% Generated human-review packet template.  Source records, Lean statements,
% and raw-source-to-Spec screening fields are substituted by
% scripts/review_dashboard_packet.py.
% Do not hand-edit generated packets; annotate the PDF or reconcile notes into
% the paper-local human-review log.
\documentclass[10pt]{article}
\usepackage[margin=0.43in,footskip=0.2in]{geometry}
\usepackage{fontspec}
% Use a Unicode-complete main face as well as a Unicode-complete code face.
% Reviewer reasons and source labels can contain exact Greek symbols outside
% verbatim blocks; silently dropping one would corrupt the review packet.
\setmainfont{DejaVu Serif}
% The broad DejaVu Sans Unicode coverage preserves Lean's mathematical glyphs
% (including U+2216 set difference and superscript identifiers) in verbatim
% review blocks.  Its proportional metrics are preferable to silently missing
% exact semantic text from a narrower monospace font.
\setmonofont{DejaVu Sans}
\usepackage{hyperref}
\usepackage{amssymb}
\usepackage{fvextra}
\usepackage{enumitem}
\setlength{\parindent}{0pt}
\setlength{\parskip}{0.15em}
\linespread{0.92}
\hypersetup{colorlinks=true,linkcolor=blue,urlcolor=blue,breaklinks=true}
\DefineVerbatimEnvironment{ReviewVerbatim}{Verbatim}{breaklines=true,breakanywhere=true,fontsize=\scriptsize}
\newcommand{\reviewerbox}{%
  \par\noindent\fbox{\parbox{0.96\linewidth}{\rule{0pt}{0.38in}}}\par
}
\newcommand{\reviewmatch}[1]{%
  \par\noindent\textbf{Reviewer check:} $\square$\; #1\par
  \noindent\emph{If left unchecked, explain the discrepancy or uncertainty below.}\par
}

\begin{document}
\fontsize{9}{10}\selectfont

\begin{center}
{\LARGE Human review packet: GGRS26CombattingGerrymanderingRCV}\\[0.35em]
\small Generated from byte-pinned source excerpts, the expanded PaperInterface
specifications, and hash-bound raw-source-to-Spec screening records.
\end{center}

\noindent\textbf{Paper:} GGRS26CombattingGerrymanderingRCV\\
\textbf{Paper id:} \texttt{GGRS26CombattingGerrymanderingRCV}\\
\textbf{Source version:} \parbox[t]{0.76\linewidth}{\raggedright arXiv:2107.07083 source bundle and byte-pinned cited publication transcript}\\
\textbf{Source record:} \texttt{audit/paper\_statement\_map.json}\\
\textbf{Rows:} 2\\
\textbf{Generated:} 2026-09-06\\
\textbf{Source URL:} \href{https://arxiv.org/pdf/2107.07083.pdf}{official source}





\section*{How to use this packet}
For each source claim, compare the exact source excerpts with the expanded
PaperInterface specification. Mark up this PDF or fill the blank reviewer box in the TeX source;
return the annotated file for reconciliation into the paper-local human-review
log.

\section*{Open the interactive dashboard}
The interactive dashboard is an optional alternative to using this PDF; it
presents the same review surface rather than an additional review requirement.
After forking and cloning (or pulling) AppliedModelingLib, run the following from the
repository root. The cache command is needed only the first time in a new clone.
\begin{ReviewVerbatim}
cd <your-repository-checkout>
lake exe cache get
python3 scripts/review_dashboard.py --paper GGRS26CombattingGerrymanderingRCV --serve
\end{ReviewVerbatim}
Open \url{http://127.0.0.1:8765/} in a browser and keep that terminal running
while reviewing. The dashboard records saved annotations in this paper's local
reviewer-owned review record.

\section*{Contents}
\begin{itemize}[leftmargin=1.2em]
\item \hyperlink{governing-declaration-section-5ef5ef0364b6939c}{Governing Lean declarations (32; supporting context)}
\item \hyperlink{paper-prerequisite-section-5ef5ef0364b6939c}{Paper-specific semantic prerequisites (24; not paper claims)}
\begin{itemize}[leftmargin=1.2em]
\item \hyperlink{paper-prerequisite-9acd2440ce1052db}{GGRS26\allowbreak{}Combatting\allowbreak{}Gerrymandering\allowbreak{}RCV.\allowbreak{}Paper\allowbreak{}Source\allowbreak{}STV\allowbreak{}State}
\item \hyperlink{paper-prerequisite-37b91baa3b4ddc9d}{GGRS26\allowbreak{}Combatting\allowbreak{}Gerrymandering\allowbreak{}RCV.\allowbreak{}Paper\allowbreak{}Thiele\allowbreak{}Weights}
\item \hyperlink{paper-prerequisite-b12522706ad3f624}{GGRS26\allowbreak{}Combatting\allowbreak{}Gerrymandering\allowbreak{}RCV.\allowbreak{}paper\_\allowbreak{}candidate\_\allowbreak{}ranked\_\allowbreak{}before}
\item \hyperlink{paper-prerequisite-a48f905958044eb5}{GGRS26\allowbreak{}Combatting\allowbreak{}Gerrymandering\allowbreak{}RCV.\allowbreak{}paper\_\allowbreak{}complete\_\allowbreak{}solid\_\allowbreak{}coalition\_\allowbreak{}rankings}
\item \hyperlink{paper-prerequisite-e34992194a8b7f0d}{GGRS26\allowbreak{}Combatting\allowbreak{}Gerrymandering\allowbreak{}RCV.\allowbreak{}paper\_\allowbreak{}first\_\allowbreak{}active\_\allowbreak{}candidate}
\item \hyperlink{paper-prerequisite-7086a653b0f05ea7}{GGRS26\allowbreak{}Combatting\allowbreak{}Gerrymandering\allowbreak{}RCV.\allowbreak{}paper\_\allowbreak{}thiele\_\allowbreak{}committee\_\allowbreak{}score}
\item \hyperlink{paper-prerequisite-517db5abc400efba}{GGRS26\allowbreak{}Combatting\allowbreak{}Gerrymandering\allowbreak{}RCV.\allowbreak{}paper\_\allowbreak{}thiele\_\allowbreak{}fixed\_\allowbreak{}size\_\allowbreak{}selected\_\allowbreak{}committee}
\item \hyperlink{paper-prerequisite-82e254229a66cc0e}{GGRS26\allowbreak{}Combatting\allowbreak{}Gerrymandering\allowbreak{}RCV.\allowbreak{}paper\_\allowbreak{}pav\_\allowbreak{}fixed\_\allowbreak{}size\_\allowbreak{}selected\_\allowbreak{}committee}
\item \hyperlink{paper-prerequisite-58801459d14912d4}{GGRS26\allowbreak{}Combatting\allowbreak{}Gerrymandering\allowbreak{}RCV.\allowbreak{}paper\_\allowbreak{}thiele\_\allowbreak{}party\_\allowbreak{}min\_\allowbreak{}argmax}
\item \hyperlink{paper-prerequisite-edccbf02c4bd6a6e}{GGRS26\allowbreak{}Combatting\allowbreak{}Gerrymandering\allowbreak{}RCV.\allowbreak{}paper\_\allowbreak{}thiele\_\allowbreak{}party\_\allowbreak{}republican\_\allowbreak{}selection}
\item \hyperlink{paper-prerequisite-5fa565502cdaaed9}{GGRS26\allowbreak{}Combatting\allowbreak{}Gerrymandering\allowbreak{}RCV.\allowbreak{}paper\_\allowbreak{}thiele\_\allowbreak{}party\_\allowbreak{}republican\_\allowbreak{}seat\_\allowbreak{}count}
\item \hyperlink{paper-prerequisite-c0f468d3ce4de7e0}{GGRS26\allowbreak{}Combatting\allowbreak{}Gerrymandering\allowbreak{}RCV.\allowbreak{}paper\_\allowbreak{}pav\_\allowbreak{}selected\_\allowbreak{}seat\_\allowbreak{}count}
\item \hyperlink{paper-prerequisite-652a33f58da6d128}{GGRS26\allowbreak{}Combatting\allowbreak{}Gerrymandering\allowbreak{}RCV.\allowbreak{}paper\_\allowbreak{}pav\_\allowbreak{}selected\_\allowbreak{}seat\_\allowbreak{}count\_\allowbreak{}spec}
\item \hyperlink{paper-prerequisite-42c2df760fca2c30}{GGRS26\allowbreak{}Combatting\allowbreak{}Gerrymandering\allowbreak{}RCV.\allowbreak{}paper\_\allowbreak{}source\_\allowbreak{}stv\_\allowbreak{}execution\_\allowbreak{}terminal}
\item \hyperlink{paper-prerequisite-587c1da71549064c}{GGRS26\allowbreak{}Combatting\allowbreak{}Gerrymandering\allowbreak{}RCV.\allowbreak{}paper\_\allowbreak{}source\_\allowbreak{}stv\_\allowbreak{}tally}
\item \hyperlink{paper-prerequisite-01a70a9b3aa00da3}{GGRS26\allowbreak{}Combatting\allowbreak{}Gerrymandering\allowbreak{}RCV.\allowbreak{}paper\_\allowbreak{}source\_\allowbreak{}stv\_\allowbreak{}d\_\allowbreak{}favoring\_\allowbreak{}transition}
\item \hyperlink{paper-prerequisite-abed5fd2c2d17fba}{GGRS26\allowbreak{}Combatting\allowbreak{}Gerrymandering\allowbreak{}RCV.\allowbreak{}paper\_\allowbreak{}source\_\allowbreak{}surplus\_\allowbreak{}transfer\_\allowbreak{}policy}
\item \hyperlink{paper-prerequisite-78828233cae137c4}{GGRS26\allowbreak{}Combatting\allowbreak{}Gerrymandering\allowbreak{}RCV.\allowbreak{}paper\_\allowbreak{}source\_\allowbreak{}stv\_\allowbreak{}rule\_\allowbreak{}execution}
\item \hyperlink{paper-prerequisite-f7db6c9d33a2a094}{GGRS26\allowbreak{}Combatting\allowbreak{}Gerrymandering\allowbreak{}RCV.\allowbreak{}paper\_\allowbreak{}thiele\_\allowbreak{}party\_\allowbreak{}democratic\_\allowbreak{}seat\_\allowbreak{}count}
\item \hyperlink{paper-prerequisite-9ae92c4421f50a56}{GGRS26\allowbreak{}Combatting\allowbreak{}Gerrymandering\allowbreak{}RCV.\allowbreak{}paper\_\allowbreak{}thiele\_\allowbreak{}party\_\allowbreak{}republican\_\allowbreak{}seat\_\allowbreak{}count\_\allowbreak{}spec}
\item \hyperlink{paper-prerequisite-1ce27606523bc987}{GGRS26\allowbreak{}Combatting\allowbreak{}Gerrymandering\allowbreak{}RCV.\allowbreak{}paper\_\allowbreak{}thiele\_\allowbreak{}party\_\allowbreak{}selected\_\allowbreak{}seat\_\allowbreak{}counts}
\item \hyperlink{paper-prerequisite-52e1d78ea1832965}{GGRS26\allowbreak{}Combatting\allowbreak{}Gerrymandering\allowbreak{}RCV.\allowbreak{}paper\_\allowbreak{}thiele\_\allowbreak{}party\_\allowbreak{}selected\_\allowbreak{}seat\_\allowbreak{}counts\_\allowbreak{}add}
\item \hyperlink{paper-prerequisite-2324cfa92955386c}{GGRS26\allowbreak{}Combatting\allowbreak{}Gerrymandering\allowbreak{}RCV.\allowbreak{}paper\_\allowbreak{}thiele\_\allowbreak{}party\_\allowbreak{}selected\_\allowbreak{}seat\_\allowbreak{}counts\_\allowbreak{}spec}
\item \hyperlink{paper-prerequisite-aea2ad46d27b40e1}{GGRS26\allowbreak{}Combatting\allowbreak{}Gerrymandering\allowbreak{}RCV.\allowbreak{}paper\_\allowbreak{}two\_\allowbreak{}party\_\allowbreak{}voter\_\allowbreak{}model}
\end{itemize}
\item \hyperlink{source-section-f10b5f68e4acf3dc}{Source claims (2)}
\begin{itemize}[leftmargin=1.2em]
\item \hyperlink{source-claim-6ccc405fd526f01f}{paper\_\allowbreak{}lemma\_\allowbreak{}c1\_\allowbreak{}pav\_\allowbreak{}selector\_\allowbreak{}eq\_\allowbreak{}unique\_\allowbreak{}integer\_\allowbreak{}interval\allowbreak{}Spec}
\item \hyperlink{source-claim-1d81d5a90a327c0b}{paper\_\allowbreak{}proposition1\_\allowbreak{}source\_\allowbreak{}selected\_\allowbreak{}stv\_\allowbreak{}and\_\allowbreak{}pav\allowbreak{}Spec}
\end{itemize}
\end{itemize}

\clearpage
\hypertarget{governing-declaration-section-5ef5ef0364b6939c}{}
\section*{Governing Lean declarations}
These exact graph-bound declaration bodies are retained by the displayed specifications or reviewed prerequisites. They are supporting context and do not add source-result rows or semantic judgments.
\hypertarget{governing-declaration-e1905fd3de83673d}{}
\subsection*{\small\ttfamily\raggedright Applied\allowbreak{}Modeling\allowbreak{}Lib.\allowbreak{}Social\allowbreak{}Choice.\allowbreak{}Voting.\allowbreak{}Approval\allowbreak{}Ballot}
\textbf{Declaration location:} reusable library
\paragraph{Lean declaration body}
\begin{ReviewVerbatim}
type:
Type u_1 → Type u_1

value:
fun (Candidate : Type u_1) => Finset Candidate
\end{ReviewVerbatim}


\hypertarget{governing-declaration-4c0320b08e4fb484}{}
\subsection*{\small\ttfamily\raggedright Applied\allowbreak{}Modeling\allowbreak{}Lib.\allowbreak{}Social\allowbreak{}Choice.\allowbreak{}Voting.\allowbreak{}Ballot}
\textbf{Declaration location:} reusable library
\paragraph{Lean declaration body}
\begin{ReviewVerbatim}
type:
Type u_1 → Type u_1

value:
fun (Candidate : Type u_1) => List Candidate
\end{ReviewVerbatim}


\hypertarget{governing-declaration-d09ca390012680b9}{}
\subsection*{\small\ttfamily\raggedright Applied\allowbreak{}Modeling\allowbreak{}Lib.\allowbreak{}Social\allowbreak{}Choice.\allowbreak{}Voting.\allowbreak{}Ballot.\allowbreak{}Valid}
\textbf{Declaration location:} reusable library
\paragraph{Lean declaration body}
\begin{ReviewVerbatim}
type:
{Candidate : Type u_1} → AppliedModelingLib.SocialChoice.Voting.Ballot Candidate → Prop

value:
fun {Candidate : Type u_1} (ballot : AppliedModelingLib.SocialChoice.Voting.Ballot Candidate) => List.Nodup ballot
\end{ReviewVerbatim}

\paragraph{Direct retained dependencies}
\begin{itemize}\raggedright
\item \hyperlink{governing-declaration-4c0320b08e4fb484}{Applied\allowbreak{}Modeling\allowbreak{}Lib.\allowbreak{}Social\allowbreak{}Choice.\allowbreak{}Voting.\allowbreak{}Ballot}
\end{itemize}
\hypertarget{governing-declaration-e1f9416a9e057b38}{}
\subsection*{\small\ttfamily\raggedright Applied\allowbreak{}Modeling\allowbreak{}Lib.\allowbreak{}Social\allowbreak{}Choice.\allowbreak{}Voting.\allowbreak{}Ballot.\allowbreak{}next\allowbreak{}Active}
\textbf{Declaration location:} reusable library
\paragraph{Lean declaration body}
\begin{ReviewVerbatim}
type:
{Candidate : Type u_1} →
  [DecidableEq Candidate] →
    AppliedModelingLib.SocialChoice.Voting.Ballot Candidate → Finset Candidate → Option Candidate

value:
fun {Candidate : Type u_1} [DecidableEq Candidate] (ballot : AppliedModelingLib.SocialChoice.Voting.Ballot Candidate)
    (active : Finset Candidate) =>
  List.brecOn ballot (AppliedModelingLib.SocialChoice.Voting.Ballot.nextActive._f active)
\end{ReviewVerbatim}

\paragraph{Direct retained dependencies}
\begin{itemize}\raggedright
\item \hyperlink{governing-declaration-4c0320b08e4fb484}{Applied\allowbreak{}Modeling\allowbreak{}Lib.\allowbreak{}Social\allowbreak{}Choice.\allowbreak{}Voting.\allowbreak{}Ballot}
\end{itemize}
\hypertarget{governing-declaration-0946df2e282443c1}{}
\subsection*{\small\ttfamily\raggedright Applied\allowbreak{}Modeling\allowbreak{}Lib.\allowbreak{}Social\allowbreak{}Choice.\allowbreak{}Voting.\allowbreak{}Ballot.\allowbreak{}active\allowbreak{}Support}
\textbf{Declaration location:} reusable library
\paragraph{Lean declaration body}
\begin{ReviewVerbatim}
type:
{Voter : Type u_1} →
  {Candidate : Type u_2} →
    [DecidableEq Candidate] →
      Finset Voter →
        (Voter → AppliedModelingLib.SocialChoice.Voting.Ballot Candidate) → Finset Candidate → Candidate → Finset Voter

value:
fun {Voter : Type u_1} {Candidate : Type u_2} [DecidableEq Candidate] (voters : Finset Voter)
    (ballots : Voter → AppliedModelingLib.SocialChoice.Voting.Ballot Candidate) (active : Finset Candidate)
    (candidate : Candidate) =>
  {voter ∈ voters | (ballots voter).nextActive active = some candidate}
\end{ReviewVerbatim}

\paragraph{Direct retained dependencies}
\begin{itemize}\raggedright
\item \hyperlink{governing-declaration-4c0320b08e4fb484}{Applied\allowbreak{}Modeling\allowbreak{}Lib.\allowbreak{}Social\allowbreak{}Choice.\allowbreak{}Voting.\allowbreak{}Ballot}
\item \hyperlink{governing-declaration-e1f9416a9e057b38}{Applied\allowbreak{}Modeling\allowbreak{}Lib.\allowbreak{}Social\allowbreak{}Choice.\allowbreak{}Voting.\allowbreak{}Ballot.\allowbreak{}next\allowbreak{}Active}
\end{itemize}
\hypertarget{governing-declaration-d2947cd1d156937d}{}
\subsection*{\small\ttfamily\raggedright Applied\allowbreak{}Modeling\allowbreak{}Lib.\allowbreak{}Social\allowbreak{}Choice.\allowbreak{}Voting.\allowbreak{}Is\allowbreak{}Min\allowbreak{}Argmax\allowbreak{}On}
\textbf{Declaration location:} reusable library
\paragraph{Lean declaration body}
\begin{ReviewVerbatim}
type:
(ℕ → ℝ) → ℕ → ℕ → Prop

value:
fun (score : ℕ → ℝ) (choice upper : ℕ) =>
  choice ≤ upper ∧
    (∀ candidate ≤ upper, score candidate ≤ score choice) ∧
      ∀ candidate ≤ upper, score candidate = score choice → choice ≤ candidate
\end{ReviewVerbatim}


\hypertarget{governing-declaration-d335879a7207e841}{}
\subsection*{\small\ttfamily\raggedright Applied\allowbreak{}Modeling\allowbreak{}Lib.\allowbreak{}Social\allowbreak{}Choice.\allowbreak{}Voting.\allowbreak{}STV\allowbreak{}Quota}
\textbf{Declaration location:} reusable library
\paragraph{Lean declaration body}
\begin{ReviewVerbatim}
type:
ℕ → ℕ → ℕ

value:
fun (seats voters : ℕ) => voters / (seats + 1) + 1
\end{ReviewVerbatim}


\hypertarget{governing-declaration-29a9f7b3ffbf6002}{}
\subsection*{\small\ttfamily\raggedright Applied\allowbreak{}Modeling\allowbreak{}Lib.\allowbreak{}Social\allowbreak{}Choice.\allowbreak{}Voting.\allowbreak{}approved\allowbreak{}Committee\allowbreak{}Count}
\textbf{Declaration location:} reusable library
\paragraph{Lean declaration body}
\begin{ReviewVerbatim}
type:
{Candidate : Type u_1} → [DecidableEq Candidate] → Finset Candidate → Finset Candidate → ℕ

value:
fun {Candidate : Type u_1} [DecidableEq Candidate] (committee ballot : Finset Candidate) => (committee ∩ ballot).card
\end{ReviewVerbatim}


\hypertarget{governing-declaration-f1687c008b580b02}{}
\subsection*{\small\ttfamily\raggedright Applied\allowbreak{}Modeling\allowbreak{}Lib.\allowbreak{}Social\allowbreak{}Choice.\allowbreak{}Voting.\allowbreak{}pav\allowbreak{}Weight}
\textbf{Declaration location:} reusable library
\paragraph{Lean declaration body}
\begin{ReviewVerbatim}
type:
ℕ → ℝ

value:
fun (approved : ℕ) => if approved = 0 then 0 else (↑approved)⁻¹
\end{ReviewVerbatim}


\hypertarget{governing-declaration-319dd7372e1f661f}{}
\subsection*{\small\ttfamily\raggedright Applied\allowbreak{}Modeling\allowbreak{}Lib.\allowbreak{}Social\allowbreak{}Choice.\allowbreak{}Voting.\allowbreak{}pav\allowbreak{}Harmonic\allowbreak{}Sum}
\textbf{Declaration location:} reusable library
\paragraph{Lean declaration body}
\begin{ReviewVerbatim}
type:
ℕ → ℝ

value:
fun (a : ℕ) => Nat.brecOn a AppliedModelingLib.SocialChoice.Voting.pavHarmonicSum._f
\end{ReviewVerbatim}

\paragraph{Direct retained dependencies}
\begin{itemize}\raggedright
\item \hyperlink{governing-declaration-f1687c008b580b02}{Applied\allowbreak{}Modeling\allowbreak{}Lib.\allowbreak{}Social\allowbreak{}Choice.\allowbreak{}Voting.\allowbreak{}pav\allowbreak{}Weight}
\end{itemize}
\hypertarget{governing-declaration-2d4e6f0580886d40}{}
\subsection*{\small\ttfamily\raggedright Applied\allowbreak{}Modeling\allowbreak{}Lib.\allowbreak{}Social\allowbreak{}Choice.\allowbreak{}Voting.\allowbreak{}rounded\allowbreak{}Seat\allowbreak{}Share}
\textbf{Declaration location:} reusable library
\paragraph{Lean declaration body}
\begin{ReviewVerbatim}
type:
ℕ → ℝ → ℕ → Prop

value:
fun (seatCount : ℕ) (partyShare : ℝ) (seats : ℕ) =>
  seatCount = ⌊partyShare * ↑seats⌋₊ ∨ seatCount = ⌈partyShare * ↑seats⌉₊
\end{ReviewVerbatim}


\hypertarget{governing-declaration-6f573ea35df34e74}{}
\subsection*{\small\ttfamily\raggedright Applied\allowbreak{}Modeling\allowbreak{}Lib.\allowbreak{}Social\allowbreak{}Choice.\allowbreak{}Voting.\allowbreak{}thiele\allowbreak{}Score}
\textbf{Declaration location:} reusable library
\paragraph{Lean declaration body}
\begin{ReviewVerbatim}
type:
{Candidate : Type u_1} →
  {Score : Type u_2} →
    [DecidableEq Candidate] →
      [AddMonoid Score] →
        (ℕ → Score) → Finset Candidate → List (AppliedModelingLib.SocialChoice.Voting.ApprovalBallot Candidate) → Score

value:
fun {Candidate : Type u_1} {Score : Type u_2} [DecidableEq Candidate] [AddMonoid Score] (weight : ℕ → Score)
    (committee : Finset Candidate) (profile : List (AppliedModelingLib.SocialChoice.Voting.ApprovalBallot Candidate)) =>
  (List.map
      (fun (ballot : Finset Candidate) =>
        weight (AppliedModelingLib.SocialChoice.Voting.approvedCommitteeCount committee ballot))
      profile).sum
\end{ReviewVerbatim}

\paragraph{Direct retained dependencies}
\begin{itemize}\raggedright
\item \hyperlink{governing-declaration-e1905fd3de83673d}{Applied\allowbreak{}Modeling\allowbreak{}Lib.\allowbreak{}Social\allowbreak{}Choice.\allowbreak{}Voting.\allowbreak{}Approval\allowbreak{}Ballot}
\item \hyperlink{governing-declaration-29a9f7b3ffbf6002}{Applied\allowbreak{}Modeling\allowbreak{}Lib.\allowbreak{}Social\allowbreak{}Choice.\allowbreak{}Voting.\allowbreak{}approved\allowbreak{}Committee\allowbreak{}Count}
\end{itemize}
\hypertarget{governing-declaration-8b872df78d871e1b}{}
\subsection*{\small\ttfamily\raggedright GGRS26\allowbreak{}Combatting\allowbreak{}Gerrymandering\allowbreak{}RCV.\allowbreak{}Ballot\allowbreak{}Routed\allowbreak{}STV\allowbreak{}State}
\textbf{Declaration location:} paper-local
\paragraph{Lean declaration body}
\begin{ReviewVerbatim}
field active:
{Voter : Type u_1} →
  {Candidate : Type u_2} →
    [inst : DecidableEq Voter] →
      [inst_1 : DecidableEq Candidate] →
        {voters : Finset Voter} →
          {initialCandidates : Finset Candidate} →
            GGRS26CombattingGerrymanderingRCV.BallotRoutedSTVState voters initialCandidates → Finset Candidate

field elected:
{Voter : Type u_1} →
  {Candidate : Type u_2} →
    [inst : DecidableEq Voter] →
      [inst_1 : DecidableEq Candidate] →
        {voters : Finset Voter} →
          {initialCandidates : Finset Candidate} →
            GGRS26CombattingGerrymanderingRCV.BallotRoutedSTVState voters initialCandidates → Finset Candidate

field weight:
{Voter : Type u_1} →
  {Candidate : Type u_2} →
    [inst : DecidableEq Voter] →
      [inst_1 : DecidableEq Candidate] →
        {voters : Finset Voter} →
          {initialCandidates : Finset Candidate} →
            GGRS26CombattingGerrymanderingRCV.BallotRoutedSTVState voters initialCandidates → Voter → ℝ

field active_subset_initial:
∀ {Voter : Type u_1} {Candidate : Type u_2} [inst : DecidableEq Voter] [inst_1 : DecidableEq Candidate]
  {voters : Finset Voter} {initialCandidates : Finset Candidate}
  (self : GGRS26CombattingGerrymanderingRCV.BallotRoutedSTVState voters initialCandidates),
  self.active ⊆ initialCandidates

field elected_subset_initial:
∀ {Voter : Type u_1} {Candidate : Type u_2} [inst : DecidableEq Voter] [inst_1 : DecidableEq Candidate]
  {voters : Finset Voter} {initialCandidates : Finset Candidate}
  (self : GGRS26CombattingGerrymanderingRCV.BallotRoutedSTVState voters initialCandidates),
  self.elected ⊆ initialCandidates

field active_elected_disjoint:
∀ {Voter : Type u_1} {Candidate : Type u_2} [inst : DecidableEq Voter] [inst_1 : DecidableEq Candidate]
  {voters : Finset Voter} {initialCandidates : Finset Candidate}
  (self : GGRS26CombattingGerrymanderingRCV.BallotRoutedSTVState voters initialCandidates),
  Disjoint self.active self.elected

field weight_nonneg:
∀ {Voter : Type u_1} {Candidate : Type u_2} [inst : DecidableEq Voter] [inst_1 : DecidableEq Candidate]
  {voters : Finset Voter} {initialCandidates : Finset Candidate}
  (self : GGRS26CombattingGerrymanderingRCV.BallotRoutedSTVState voters initialCandidates),
  ∀ voter ∈ voters, 0 ≤ self.weight voter
\end{ReviewVerbatim}


\hypertarget{governing-declaration-e4f70d25375921c7}{}
\subsection*{\small\ttfamily\raggedright GGRS26\allowbreak{}Combatting\allowbreak{}Gerrymandering\allowbreak{}RCV.\allowbreak{}Ballot\allowbreak{}Routed\allowbreak{}STV\allowbreak{}State.\allowbreak{}initial}
\textbf{Declaration location:} paper-local
\paragraph{Lean declaration body}
\begin{ReviewVerbatim}
type:
{Voter : Type u_1} →
  {Candidate : Type u_2} →
    [inst : DecidableEq Voter] →
      [inst_1 : DecidableEq Candidate] →
        {voters : Finset Voter} →
          {initialCandidates : Finset Candidate} →
            (initialWeight : Voter → ℝ) →
              (∀ voter ∈ voters, 0 ≤ initialWeight voter) →
                GGRS26CombattingGerrymanderingRCV.BallotRoutedSTVState voters initialCandidates

value:
fun {Voter : Type u_1} {Candidate : Type u_2} [DecidableEq Voter] [DecidableEq Candidate] {voters : Finset Voter}
    {initialCandidates : Finset Candidate} (initialWeight : Voter → ℝ)
    (hinitialWeight_nonneg : ∀ voter ∈ voters, 0 ≤ initialWeight voter) =>
  { active := initialCandidates, elected := ∅, weight := initialWeight,
    active_subset_initial := fun {{candidate : Candidate}} (hcandidate : candidate ∈ initialCandidates) => hcandidate,
    elected_subset_initial := GGRS26CombattingGerrymanderingRCV.BallotRoutedSTVState.initial._proof_1,
    active_elected_disjoint := GGRS26CombattingGerrymanderingRCV.BallotRoutedSTVState.initial._proof_2,
    weight_nonneg := hinitialWeight_nonneg }
\end{ReviewVerbatim}

\paragraph{Direct retained dependencies}
\begin{itemize}\raggedright
\item \hyperlink{governing-declaration-8b872df78d871e1b}{GGRS26\allowbreak{}Combatting\allowbreak{}Gerrymandering\allowbreak{}RCV.\allowbreak{}Ballot\allowbreak{}Routed\allowbreak{}STV\allowbreak{}State}
\end{itemize}
\hypertarget{governing-declaration-27b5b75da6f75108}{}
\subsection*{\small\ttfamily\raggedright GGRS26\allowbreak{}Combatting\allowbreak{}Gerrymandering\allowbreak{}RCV.\allowbreak{}Ballot\allowbreak{}Routed\allowbreak{}STV\allowbreak{}Terminal}
\textbf{Declaration location:} paper-local
\paragraph{Lean declaration body}
\begin{ReviewVerbatim}
type:
{Voter : Type u_1} →
  {Candidate : Type u_2} →
    [inst : DecidableEq Voter] →
      [inst_1 : DecidableEq Candidate] →
        {voters : Finset Voter} →
          {initialCandidates : Finset Candidate} →
            ℕ → GGRS26CombattingGerrymanderingRCV.BallotRoutedSTVState voters initialCandidates → Prop

value:
fun {Voter : Type u_1} {Candidate : Type u_2} [DecidableEq Voter] [DecidableEq Candidate] {voters : Finset Voter}
    {initialCandidates : Finset Candidate} (seats : ℕ)
    (state : GGRS26CombattingGerrymanderingRCV.BallotRoutedSTVState voters initialCandidates) =>
  state.elected.card = seats ∨ state.active.card = seats - state.elected.card
\end{ReviewVerbatim}

\paragraph{Direct retained dependencies}
\begin{itemize}\raggedright
\item \hyperlink{governing-declaration-8b872df78d871e1b}{GGRS26\allowbreak{}Combatting\allowbreak{}Gerrymandering\allowbreak{}RCV.\allowbreak{}Ballot\allowbreak{}Routed\allowbreak{}STV\allowbreak{}State}
\end{itemize}
\hypertarget{governing-declaration-3adeede1bc778473}{}
\subsection*{\small\ttfamily\raggedright GGRS26\allowbreak{}Combatting\allowbreak{}Gerrymandering\allowbreak{}RCV.\allowbreak{}Party\allowbreak{}Approval\allowbreak{}Ballot}
\textbf{Declaration location:} paper-local
\paragraph{Lean declaration body}
\begin{ReviewVerbatim}
type:
Type u_1 → Type u_1

value:
fun (Candidate : Type u_1) => AppliedModelingLib.SocialChoice.Voting.ApprovalBallot Candidate
\end{ReviewVerbatim}

\paragraph{Direct retained dependencies}
\begin{itemize}\raggedright
\item \hyperlink{governing-declaration-e1905fd3de83673d}{Applied\allowbreak{}Modeling\allowbreak{}Lib.\allowbreak{}Social\allowbreak{}Choice.\allowbreak{}Voting.\allowbreak{}Approval\allowbreak{}Ballot}
\end{itemize}
\hypertarget{governing-declaration-21463db3a6da8b5e}{}
\subsection*{\small\ttfamily\raggedright GGRS26\allowbreak{}Combatting\allowbreak{}Gerrymandering\allowbreak{}RCV.\allowbreak{}ballot\allowbreak{}Routed\allowbreak{}Party\allowbreak{}Final\allowbreak{}Seats}
\textbf{Declaration location:} paper-local
\paragraph{Lean declaration body}
\begin{ReviewVerbatim}
type:
{Voter : Type u_1} →
  {Candidate : Type u_2} →
    [inst : DecidableEq Voter] →
      [inst_1 : DecidableEq Candidate] →
        {voters : Finset Voter} →
          {initialCandidates : Finset Candidate} →
            Finset Candidate → ℕ → GGRS26CombattingGerrymanderingRCV.BallotRoutedSTVState voters initialCandidates → ℕ

value:
fun {Voter : Type u_1} {Candidate : Type u_2} [DecidableEq Voter] [DecidableEq Candidate] {voters : Finset Voter}
    {initialCandidates : Finset Candidate} (partyCandidates : Finset Candidate) (seats : ℕ)
    (state : GGRS26CombattingGerrymanderingRCV.BallotRoutedSTVState voters initialCandidates) =>
  (state.elected ∩ partyCandidates).card +
    if state.elected.card < seats then (state.active ∩ partyCandidates).card else 0
\end{ReviewVerbatim}

\paragraph{Direct retained dependencies}
\begin{itemize}\raggedright
\item \hyperlink{governing-declaration-8b872df78d871e1b}{GGRS26\allowbreak{}Combatting\allowbreak{}Gerrymandering\allowbreak{}RCV.\allowbreak{}Ballot\allowbreak{}Routed\allowbreak{}STV\allowbreak{}State}
\end{itemize}
\hypertarget{governing-declaration-e25badf495bb0753}{}
\subsection*{\small\ttfamily\raggedright GGRS26\allowbreak{}Combatting\allowbreak{}Gerrymandering\allowbreak{}RCV.\allowbreak{}ballot\allowbreak{}Routed\allowbreak{}Support}
\textbf{Declaration location:} paper-local
\paragraph{Lean declaration body}
\begin{ReviewVerbatim}
type:
{Voter : Type u_1} →
  {Candidate : Type u_2} →
    [DecidableEq Candidate] →
      Finset Voter →
        (Voter → AppliedModelingLib.SocialChoice.Voting.Ballot Candidate) → Finset Candidate → Candidate → Finset Voter

value:
fun {Voter : Type u_1} {Candidate : Type u_2} [DecidableEq Candidate] (voters : Finset Voter)
    (ballots : Voter → AppliedModelingLib.SocialChoice.Voting.Ballot Candidate) (active : Finset Candidate)
    (candidate : Candidate) =>
  AppliedModelingLib.SocialChoice.Voting.Ballot.activeSupport voters ballots active candidate
\end{ReviewVerbatim}

\paragraph{Direct retained dependencies}
\begin{itemize}\raggedright
\item \hyperlink{governing-declaration-4c0320b08e4fb484}{Applied\allowbreak{}Modeling\allowbreak{}Lib.\allowbreak{}Social\allowbreak{}Choice.\allowbreak{}Voting.\allowbreak{}Ballot}
\item \hyperlink{governing-declaration-0946df2e282443c1}{Applied\allowbreak{}Modeling\allowbreak{}Lib.\allowbreak{}Social\allowbreak{}Choice.\allowbreak{}Voting.\allowbreak{}Ballot.\allowbreak{}active\allowbreak{}Support}
\end{itemize}
\hypertarget{governing-declaration-9546927b87891fe9}{}
\subsection*{\small\ttfamily\raggedright GGRS26\allowbreak{}Combatting\allowbreak{}Gerrymandering\allowbreak{}RCV.\allowbreak{}ballot\allowbreak{}Routed\allowbreak{}Tally}
\textbf{Declaration location:} paper-local
\paragraph{Lean declaration body}
\begin{ReviewVerbatim}
type:
{Voter : Type u_1} →
  {Candidate : Type u_2} →
    [DecidableEq Candidate] →
      Finset Voter →
        (Voter → AppliedModelingLib.SocialChoice.Voting.Ballot Candidate) →
          Finset Candidate → (Voter → ℝ) → Candidate → ℝ

value:
fun {Voter : Type u_1} {Candidate : Type u_2} [DecidableEq Candidate] (voters : Finset Voter)
    (ballots : Voter → AppliedModelingLib.SocialChoice.Voting.Ballot Candidate) (active : Finset Candidate)
    (weight : Voter → ℝ) (candidate : Candidate) =>
  ∑ voter ∈ GGRS26CombattingGerrymanderingRCV.ballotRoutedSupport voters ballots active candidate, weight voter
\end{ReviewVerbatim}

\paragraph{Direct retained dependencies}
\begin{itemize}\raggedright
\item \hyperlink{governing-declaration-4c0320b08e4fb484}{Applied\allowbreak{}Modeling\allowbreak{}Lib.\allowbreak{}Social\allowbreak{}Choice.\allowbreak{}Voting.\allowbreak{}Ballot}
\item \hyperlink{governing-declaration-e25badf495bb0753}{GGRS26\allowbreak{}Combatting\allowbreak{}Gerrymandering\allowbreak{}RCV.\allowbreak{}ballot\allowbreak{}Routed\allowbreak{}Support}
\end{itemize}
\hypertarget{governing-declaration-9554b7918969e197}{}
\subsection*{\small\ttfamily\raggedright GGRS26\allowbreak{}Combatting\allowbreak{}Gerrymandering\allowbreak{}RCV.\allowbreak{}Ballot\allowbreak{}Routed\allowbreak{}STV\allowbreak{}Transfer\allowbreak{}Policy}
\textbf{Declaration location:} paper-local
\paragraph{Lean declaration body}
\begin{ReviewVerbatim}
field electUpdate:
{Voter : Type u_1} →
  {Candidate : Type u_2} →
    [inst : DecidableEq Voter] →
      [inst_1 : DecidableEq Candidate] →
        {voters : Finset Voter} →
          {ballots : Voter → AppliedModelingLib.SocialChoice.Voting.Ballot Candidate} →
            {quota : ℝ} →
              GGRS26CombattingGerrymanderingRCV.BallotRoutedSTVTransferPolicy voters ballots quota →
                Finset Candidate → Candidate → (Voter → ℝ) → (Voter → ℝ) → Prop

field eliminateUpdate:
{Voter : Type u_1} →
  {Candidate : Type u_2} →
    [inst : DecidableEq Voter] →
      [inst_1 : DecidableEq Candidate] →
        {voters : Finset Voter} →
          {ballots : Voter → AppliedModelingLib.SocialChoice.Voting.Ballot Candidate} →
            {quota : ℝ} →
              GGRS26CombattingGerrymanderingRCV.BallotRoutedSTVTransferPolicy voters ballots quota →
                Finset Candidate → Candidate → (Voter → ℝ) → (Voter → ℝ) → Prop

field elect_exists:
∀ {Voter : Type u_1} {Candidate : Type u_2} [inst : DecidableEq Voter] [inst_1 : DecidableEq Candidate]
  {voters : Finset Voter} {ballots : Voter → AppliedModelingLib.SocialChoice.Voting.Ballot Candidate} {quota : ℝ}
  (self : GGRS26CombattingGerrymanderingRCV.BallotRoutedSTVTransferPolicy voters ballots quota)
  (active : Finset Candidate) (winner : Candidate) (beforeWeight : Voter → ℝ),
  (∀ voter ∈ voters, 0 ≤ beforeWeight voter) →
    winner ∈ active →
      quota ≤ GGRS26CombattingGerrymanderingRCV.ballotRoutedTally voters ballots active beforeWeight winner →
        ∃ (afterWeight : Voter → ℝ), self.electUpdate active winner beforeWeight afterWeight

field eliminate_exists:
∀ {Voter : Type u_1} {Candidate : Type u_2} [inst : DecidableEq Voter] [inst_1 : DecidableEq Candidate]
  {voters : Finset Voter} {ballots : Voter → AppliedModelingLib.SocialChoice.Voting.Ballot Candidate} {quota : ℝ}
  (self : GGRS26CombattingGerrymanderingRCV.BallotRoutedSTVTransferPolicy voters ballots quota)
  (active : Finset Candidate) (loser : Candidate) (beforeWeight : Voter → ℝ),
  (∀ voter ∈ voters, 0 ≤ beforeWeight voter) →
    loser ∈ active → ∃ (afterWeight : Voter → ℝ), self.eliminateUpdate active loser beforeWeight afterWeight

field elect_preserves_nonneg:
∀ {Voter : Type u_1} {Candidate : Type u_2} [inst : DecidableEq Voter] [inst_1 : DecidableEq Candidate]
  {voters : Finset Voter} {ballots : Voter → AppliedModelingLib.SocialChoice.Voting.Ballot Candidate} {quota : ℝ}
  (self : GGRS26CombattingGerrymanderingRCV.BallotRoutedSTVTransferPolicy voters ballots quota)
  (active : Finset Candidate) (winner : Candidate) (beforeWeight afterWeight : Voter → ℝ),
  (∀ voter ∈ voters, 0 ≤ beforeWeight voter) →
    self.electUpdate active winner beforeWeight afterWeight → ∀ voter ∈ voters, 0 ≤ afterWeight voter

field elect_unchanged_off_support:
∀ {Voter : Type u_1} {Candidate : Type u_2} [inst : DecidableEq Voter] [inst_1 : DecidableEq Candidate]
  {voters : Finset Voter} {ballots : Voter → AppliedModelingLib.SocialChoice.Voting.Ballot Candidate} {quota : ℝ}
  (self : GGRS26CombattingGerrymanderingRCV.BallotRoutedSTVTransferPolicy voters ballots quota)
  (active : Finset Candidate) (winner : Candidate) (beforeWeight afterWeight : Voter → ℝ),
  (∀ voter ∈ voters, 0 ≤ beforeWeight voter) →
    self.electUpdate active winner beforeWeight afterWeight →
      ∀ voter ∈ voters,
        voter ∉ GGRS26CombattingGerrymanderingRCV.ballotRoutedSupport voters ballots active winner →
          afterWeight voter = beforeWeight voter

field elect_support_mass_drop_exactly_quota:
∀ {Voter : Type u_1} {Candidate : Type u_2} [inst : DecidableEq Voter] [inst_1 : DecidableEq Candidate]
  {voters : Finset Voter} {ballots : Voter → AppliedModelingLib.SocialChoice.Voting.Ballot Candidate} {quota : ℝ}
  (self : GGRS26CombattingGerrymanderingRCV.BallotRoutedSTVTransferPolicy voters ballots quota)
  (active : Finset Candidate) (winner : Candidate) (beforeWeight afterWeight : Voter → ℝ),
  (∀ voter ∈ voters, 0 ≤ beforeWeight voter) →
    self.electUpdate active winner beforeWeight afterWeight →
      ∑ voter ∈ GGRS26CombattingGerrymanderingRCV.ballotRoutedSupport voters ballots active winner, afterWeight voter =
        ∑ voter ∈ GGRS26CombattingGerrymanderingRCV.ballotRoutedSupport voters ballots active winner,
            beforeWeight voter -
          quota

field eliminate_weight_unchanged:
∀ {Voter : Type u_1} {Candidate : Type u_2} [inst : DecidableEq Voter] [inst_1 : DecidableEq Candidate]
  {voters : Finset Voter} {ballots : Voter → AppliedModelingLib.SocialChoice.Voting.Ballot Candidate} {quota : ℝ}
  (self : GGRS26CombattingGerrymanderingRCV.BallotRoutedSTVTransferPolicy voters ballots quota)
  (active : Finset Candidate) (loser : Candidate) (beforeWeight afterWeight : Voter → ℝ),
  (∀ voter ∈ voters, 0 ≤ beforeWeight voter) →
    self.eliminateUpdate active loser beforeWeight afterWeight →
      ∀ voter ∈ voters, afterWeight voter = beforeWeight voter
\end{ReviewVerbatim}

\paragraph{Direct retained dependencies}
\begin{itemize}\raggedright
\item \hyperlink{governing-declaration-4c0320b08e4fb484}{Applied\allowbreak{}Modeling\allowbreak{}Lib.\allowbreak{}Social\allowbreak{}Choice.\allowbreak{}Voting.\allowbreak{}Ballot}
\item \hyperlink{governing-declaration-e25badf495bb0753}{GGRS26\allowbreak{}Combatting\allowbreak{}Gerrymandering\allowbreak{}RCV.\allowbreak{}ballot\allowbreak{}Routed\allowbreak{}Support}
\item \hyperlink{governing-declaration-9546927b87891fe9}{GGRS26\allowbreak{}Combatting\allowbreak{}Gerrymandering\allowbreak{}RCV.\allowbreak{}ballot\allowbreak{}Routed\allowbreak{}Tally}
\end{itemize}
\hypertarget{governing-declaration-4dcbe128f2b022b3}{}
\subsection*{\small\ttfamily\raggedright GGRS26\allowbreak{}Combatting\allowbreak{}Gerrymandering\allowbreak{}RCV.\allowbreak{}Ballot\allowbreak{}Routed\allowbreak{}D\allowbreak{}Favoring\allowbreak{}STV\allowbreak{}Transition}
\textbf{Declaration location:} paper-local
\paragraph{Lean declaration body}
\begin{ReviewVerbatim}
type:
{Voter : Type u_1} →
  {Candidate : Type u_2} →
    [inst : DecidableEq Voter] →
      [inst_1 : DecidableEq Candidate] →
        {voters : Finset Voter} →
          {initialCandidates : Finset Candidate} →
            {favoredParty : Finset Candidate} →
              (ballots : Voter → AppliedModelingLib.SocialChoice.Voting.Ballot Candidate) →
                (quota : ℝ) →
                  GGRS26CombattingGerrymanderingRCV.BallotRoutedSTVTransferPolicy voters ballots quota →
                    ℕ →
                      GGRS26CombattingGerrymanderingRCV.BallotRoutedSTVState voters initialCandidates →
                        GGRS26CombattingGerrymanderingRCV.BallotRoutedSTVState voters initialCandidates → Prop

constructor GGRS26CombattingGerrymanderingRCV.BallotRoutedDFavoringSTVTransition.elect:
∀ {Voter : Type u_1} {Candidate : Type u_2} [inst : DecidableEq Voter] [inst_1 : DecidableEq Candidate]
  {voters : Finset Voter} {initialCandidates favoredParty : Finset Candidate}
  {ballots : Voter → AppliedModelingLib.SocialChoice.Voting.Ballot Candidate} {quota : ℝ}
  {policy : GGRS26CombattingGerrymanderingRCV.BallotRoutedSTVTransferPolicy voters ballots quota} {seats : ℕ}
  {before after : GGRS26CombattingGerrymanderingRCV.BallotRoutedSTVState voters initialCandidates} (winner : Candidate),
  ¬GGRS26CombattingGerrymanderingRCV.BallotRoutedSTVTerminal seats before →
    winner ∈ before.active →
      before.elected.card < seats →
        quota ≤ GGRS26CombattingGerrymanderingRCV.ballotRoutedTally voters ballots before.active before.weight winner →
          (winner ∉ favoredParty →
              ∀ candidate ∈ before.active,
                GGRS26CombattingGerrymanderingRCV.ballotRoutedTally voters ballots before.active before.weight
                      candidate =
                    GGRS26CombattingGerrymanderingRCV.ballotRoutedTally voters ballots before.active before.weight
                      winner →
                  candidate ∉ favoredParty) →
            after.active = before.active.erase winner →
              after.elected = insert winner before.elected →
                policy.electUpdate before.active winner before.weight after.weight →
                  GGRS26CombattingGerrymanderingRCV.BallotRoutedDFavoringSTVTransition ballots quota policy seats before
                    after

constructor GGRS26CombattingGerrymanderingRCV.BallotRoutedDFavoringSTVTransition.eliminate:
∀ {Voter : Type u_1} {Candidate : Type u_2} [inst : DecidableEq Voter] [inst_1 : DecidableEq Candidate]
  {voters : Finset Voter} {initialCandidates favoredParty : Finset Candidate}
  {ballots : Voter → AppliedModelingLib.SocialChoice.Voting.Ballot Candidate} {quota : ℝ}
  {policy : GGRS26CombattingGerrymanderingRCV.BallotRoutedSTVTransferPolicy voters ballots quota} {seats : ℕ}
  {before after : GGRS26CombattingGerrymanderingRCV.BallotRoutedSTVState voters initialCandidates} (loser : Candidate),
  ¬GGRS26CombattingGerrymanderingRCV.BallotRoutedSTVTerminal seats before →
    loser ∈ before.active →
      (∀ candidate ∈ before.active,
          GGRS26CombattingGerrymanderingRCV.ballotRoutedTally voters ballots before.active before.weight candidate <
            quota) →
        (∀ candidate ∈ before.active,
            GGRS26CombattingGerrymanderingRCV.ballotRoutedTally voters ballots before.active before.weight loser ≤
              GGRS26CombattingGerrymanderingRCV.ballotRoutedTally voters ballots before.active before.weight
                candidate) →
          (loser ∈ favoredParty →
              ∀ candidate ∈ before.active,
                GGRS26CombattingGerrymanderingRCV.ballotRoutedTally voters ballots before.active before.weight
                      candidate =
                    GGRS26CombattingGerrymanderingRCV.ballotRoutedTally voters ballots before.active before.weight
                      loser →
                  candidate ∈ favoredParty) →
            after.active = before.active.erase loser →
              after.elected = before.elected →
                policy.eliminateUpdate before.active loser before.weight after.weight →
                  GGRS26CombattingGerrymanderingRCV.BallotRoutedDFavoringSTVTransition ballots quota policy seats before
                    after
\end{ReviewVerbatim}

\paragraph{Direct retained dependencies}
\begin{itemize}\raggedright
\item \hyperlink{governing-declaration-4c0320b08e4fb484}{Applied\allowbreak{}Modeling\allowbreak{}Lib.\allowbreak{}Social\allowbreak{}Choice.\allowbreak{}Voting.\allowbreak{}Ballot}
\item \hyperlink{governing-declaration-8b872df78d871e1b}{GGRS26\allowbreak{}Combatting\allowbreak{}Gerrymandering\allowbreak{}RCV.\allowbreak{}Ballot\allowbreak{}Routed\allowbreak{}STV\allowbreak{}State}
\item \hyperlink{governing-declaration-27b5b75da6f75108}{GGRS26\allowbreak{}Combatting\allowbreak{}Gerrymandering\allowbreak{}RCV.\allowbreak{}Ballot\allowbreak{}Routed\allowbreak{}STV\allowbreak{}Terminal}
\item \hyperlink{governing-declaration-9554b7918969e197}{GGRS26\allowbreak{}Combatting\allowbreak{}Gerrymandering\allowbreak{}RCV.\allowbreak{}Ballot\allowbreak{}Routed\allowbreak{}STV\allowbreak{}Transfer\allowbreak{}Policy}
\item \hyperlink{governing-declaration-9546927b87891fe9}{GGRS26\allowbreak{}Combatting\allowbreak{}Gerrymandering\allowbreak{}RCV.\allowbreak{}ballot\allowbreak{}Routed\allowbreak{}Tally}
\end{itemize}
\hypertarget{governing-declaration-3bfffc20e8b6444f}{}
\subsection*{\small\ttfamily\raggedright GGRS26\allowbreak{}Combatting\allowbreak{}Gerrymandering\allowbreak{}RCV.\allowbreak{}Ballot\allowbreak{}Routed\allowbreak{}D\allowbreak{}Favoring\allowbreak{}STV\allowbreak{}Run}
\textbf{Declaration location:} paper-local
\paragraph{Lean declaration body}
\begin{ReviewVerbatim}
type:
{Voter : Type u_1} →
  {Candidate : Type u_2} →
    [inst : DecidableEq Voter] →
      [inst_1 : DecidableEq Candidate] →
        {voters : Finset Voter} →
          {initialCandidates : Finset Candidate} →
            {favoredParty : Finset Candidate} →
              (ballots : Voter → AppliedModelingLib.SocialChoice.Voting.Ballot Candidate) →
                (quota : ℝ) →
                  GGRS26CombattingGerrymanderingRCV.BallotRoutedSTVTransferPolicy voters ballots quota →
                    ℕ →
                      GGRS26CombattingGerrymanderingRCV.BallotRoutedSTVState voters initialCandidates →
                        GGRS26CombattingGerrymanderingRCV.BallotRoutedSTVState voters initialCandidates → Prop

value:
fun {Voter : Type u_1} {Candidate : Type u_2} [DecidableEq Voter] [DecidableEq Candidate] {voters : Finset Voter}
    {initialCandidates favoredParty : Finset Candidate}
    (ballots : Voter → AppliedModelingLib.SocialChoice.Voting.Ballot Candidate) (quota : ℝ)
    (policy : GGRS26CombattingGerrymanderingRCV.BallotRoutedSTVTransferPolicy voters ballots quota) (seats : ℕ)
    (initial terminal : GGRS26CombattingGerrymanderingRCV.BallotRoutedSTVState voters initialCandidates) =>
  Relation.ReflTransGen
    (GGRS26CombattingGerrymanderingRCV.BallotRoutedDFavoringSTVTransition ballots quota policy seats) initial terminal
\end{ReviewVerbatim}

\paragraph{Direct retained dependencies}
\begin{itemize}\raggedright
\item \hyperlink{governing-declaration-4c0320b08e4fb484}{Applied\allowbreak{}Modeling\allowbreak{}Lib.\allowbreak{}Social\allowbreak{}Choice.\allowbreak{}Voting.\allowbreak{}Ballot}
\item \hyperlink{governing-declaration-4dcbe128f2b022b3}{GGRS26\allowbreak{}Combatting\allowbreak{}Gerrymandering\allowbreak{}RCV.\allowbreak{}Ballot\allowbreak{}Routed\allowbreak{}D\allowbreak{}Favoring\allowbreak{}STV\allowbreak{}Transition}
\item \hyperlink{governing-declaration-8b872df78d871e1b}{GGRS26\allowbreak{}Combatting\allowbreak{}Gerrymandering\allowbreak{}RCV.\allowbreak{}Ballot\allowbreak{}Routed\allowbreak{}STV\allowbreak{}State}
\item \hyperlink{governing-declaration-9554b7918969e197}{GGRS26\allowbreak{}Combatting\allowbreak{}Gerrymandering\allowbreak{}RCV.\allowbreak{}Ballot\allowbreak{}Routed\allowbreak{}STV\allowbreak{}Transfer\allowbreak{}Policy}
\end{itemize}
\hypertarget{governing-declaration-3a02b26960fad588}{}
\subsection*{\small\ttfamily\raggedright GGRS26\allowbreak{}Combatting\allowbreak{}Gerrymandering\allowbreak{}RCV.\allowbreak{}paper\_\allowbreak{}complete\_\allowbreak{}ranked\_\allowbreak{}ballots}
\textbf{Declaration location:} paper-local
\paragraph{Lean declaration body}
\begin{ReviewVerbatim}
type:
{Voter : Type u_1} →
  {Candidate : Type u_2} →
    [DecidableEq Candidate] →
      Finset Voter → (Voter → AppliedModelingLib.SocialChoice.Voting.Ballot Candidate) → Finset Candidate → Prop

value:
fun {Voter : Type u_1} {Candidate : Type u_2} [DecidableEq Candidate] (voters : Finset Voter)
    (ballots : Voter → AppliedModelingLib.SocialChoice.Voting.Ballot Candidate) (allCandidates : Finset Candidate) =>
  ∀ voter ∈ voters,
    (ballots voter).Valid ∧ ∀ (candidate : Candidate), candidate ∈ ballots voter ↔ candidate ∈ allCandidates
\end{ReviewVerbatim}

\paragraph{Direct retained dependencies}
\begin{itemize}\raggedright
\item \hyperlink{governing-declaration-4c0320b08e4fb484}{Applied\allowbreak{}Modeling\allowbreak{}Lib.\allowbreak{}Social\allowbreak{}Choice.\allowbreak{}Voting.\allowbreak{}Ballot}
\item \hyperlink{governing-declaration-d09ca390012680b9}{Applied\allowbreak{}Modeling\allowbreak{}Lib.\allowbreak{}Social\allowbreak{}Choice.\allowbreak{}Voting.\allowbreak{}Ballot.\allowbreak{}Valid}
\end{itemize}
\hypertarget{governing-declaration-ab89c07c272f0ccf}{}
\subsection*{\small\ttfamily\raggedright GGRS26\allowbreak{}Combatting\allowbreak{}Gerrymandering\allowbreak{}RCV.\allowbreak{}paper\_\allowbreak{}complete\_\allowbreak{}party\_\allowbreak{}ranking\_\allowbreak{}ballots}
\textbf{Declaration location:} paper-local
\paragraph{Lean declaration body}
\begin{ReviewVerbatim}
type:
{Voter : Type u_1} →
  {Candidate : Type u_2} →
    [DecidableEq Candidate] →
      Finset Voter →
        (Voter → AppliedModelingLib.SocialChoice.Voting.Ballot Candidate) → Finset Candidate → Finset Candidate → Prop

value:
fun {Voter : Type u_1} {Candidate : Type u_2} [DecidableEq Candidate] (voters : Finset Voter)
    (ballots : Voter → AppliedModelingLib.SocialChoice.Voting.Ballot Candidate)
    (partyCandidates allCandidates : Finset Candidate) =>
  GGRS26CombattingGerrymanderingRCV.paper_complete_ranked_ballots voters ballots allCandidates ∧
    ∀ voter ∈ voters,
      ∀ (active : Finset Candidate),
        (∃ same ∈ partyCandidates, same ∈ active) →
          ∃ same ∈ partyCandidates, (ballots voter).nextActive active = some same
\end{ReviewVerbatim}

\paragraph{Direct retained dependencies}
\begin{itemize}\raggedright
\item \hyperlink{governing-declaration-4c0320b08e4fb484}{Applied\allowbreak{}Modeling\allowbreak{}Lib.\allowbreak{}Social\allowbreak{}Choice.\allowbreak{}Voting.\allowbreak{}Ballot}
\item \hyperlink{governing-declaration-e1f9416a9e057b38}{Applied\allowbreak{}Modeling\allowbreak{}Lib.\allowbreak{}Social\allowbreak{}Choice.\allowbreak{}Voting.\allowbreak{}Ballot.\allowbreak{}next\allowbreak{}Active}
\item \hyperlink{governing-declaration-3a02b26960fad588}{GGRS26\allowbreak{}Combatting\allowbreak{}Gerrymandering\allowbreak{}RCV.\allowbreak{}paper\_\allowbreak{}complete\_\allowbreak{}ranked\_\allowbreak{}ballots}
\end{itemize}
\hypertarget{governing-declaration-2864a4aebe1e6e50}{}
\subsection*{\small\ttfamily\raggedright GGRS26\allowbreak{}Combatting\allowbreak{}Gerrymandering\allowbreak{}RCV.\allowbreak{}paper\_\allowbreak{}pav\_\allowbreak{}thiele\_\allowbreak{}weights}
\textbf{Declaration location:} paper-local
\paragraph{Lean declaration body}
\begin{ReviewVerbatim}
type:
GGRS26CombattingGerrymanderingRCV.PaperThieleWeights

value:
{ marginalWeight := AppliedModelingLib.SocialChoice.Voting.pavWeight,
  marginal_nonincreasing := GGRS26CombattingGerrymanderingRCV.paper_pav_thiele_weights._proof_1,
  voterScore := AppliedModelingLib.SocialChoice.Voting.pavHarmonicSum,
  voterScore_zero := GGRS26CombattingGerrymanderingRCV.paper_pav_thiele_weights._proof_2,
  voterScore_succ := GGRS26CombattingGerrymanderingRCV.paper_pav_thiele_weights._proof_3 }
\end{ReviewVerbatim}

\paragraph{Direct retained dependencies}
\begin{itemize}\raggedright
\item \hyperlink{governing-declaration-319dd7372e1f661f}{Applied\allowbreak{}Modeling\allowbreak{}Lib.\allowbreak{}Social\allowbreak{}Choice.\allowbreak{}Voting.\allowbreak{}pav\allowbreak{}Harmonic\allowbreak{}Sum}
\item \hyperlink{governing-declaration-f1687c008b580b02}{Applied\allowbreak{}Modeling\allowbreak{}Lib.\allowbreak{}Social\allowbreak{}Choice.\allowbreak{}Voting.\allowbreak{}pav\allowbreak{}Weight}
\item \hyperlink{paper-prerequisite-37b91baa3b4ddc9d}{GGRS26\allowbreak{}Combatting\allowbreak{}Gerrymandering\allowbreak{}RCV.\allowbreak{}Paper\allowbreak{}Thiele\allowbreak{}Weights}
\end{itemize}
\hypertarget{governing-declaration-a26378e50367c86c}{}
\subsection*{\small\ttfamily\raggedright GGRS26\allowbreak{}Combatting\allowbreak{}Gerrymandering\allowbreak{}RCV.\allowbreak{}paper\_\allowbreak{}proposition1\_\allowbreak{}ballot\_\allowbreak{}routed\_\allowbreak{}stv\_\allowbreak{}initial\_\allowbreak{}state}
\textbf{Declaration location:} paper-local
\paragraph{Lean declaration body}
\begin{ReviewVerbatim}
type:
{Voter : Type u_1} →
  {Candidate : Type u_2} →
    [inst : DecidableEq Voter] →
      [inst_1 : DecidableEq Candidate] →
        (allVoters : Finset Voter) →
          (initialActive : Finset Candidate) →
            GGRS26CombattingGerrymanderingRCV.BallotRoutedSTVState allVoters initialActive

value:
fun {Voter : Type u_1} {Candidate : Type u_2} [DecidableEq Voter] [DecidableEq Candidate] (allVoters : Finset Voter)
    (initialActive : Finset Candidate) =>
  GGRS26CombattingGerrymanderingRCV.BallotRoutedSTVState.initial (fun (x : Voter) => 1)
    (GGRS26CombattingGerrymanderingRCV.paper_proposition1_ballot_routed_stv_initial_state._proof_1 allVoters)
\end{ReviewVerbatim}

\paragraph{Direct retained dependencies}
\begin{itemize}\raggedright
\item \hyperlink{governing-declaration-8b872df78d871e1b}{GGRS26\allowbreak{}Combatting\allowbreak{}Gerrymandering\allowbreak{}RCV.\allowbreak{}Ballot\allowbreak{}Routed\allowbreak{}STV\allowbreak{}State}
\item \hyperlink{governing-declaration-e4f70d25375921c7}{GGRS26\allowbreak{}Combatting\allowbreak{}Gerrymandering\allowbreak{}RCV.\allowbreak{}Ballot\allowbreak{}Routed\allowbreak{}STV\allowbreak{}State.\allowbreak{}initial}
\end{itemize}
\hypertarget{governing-declaration-d8487ebdee5a1786}{}
\subsection*{\small\ttfamily\raggedright GGRS26\allowbreak{}Combatting\allowbreak{}Gerrymandering\allowbreak{}RCV.\allowbreak{}paper\_\allowbreak{}thiele\_\allowbreak{}party\_\allowbreak{}seat\_\allowbreak{}score}
\textbf{Declaration location:} paper-local
\paragraph{Lean declaration body}
\begin{ReviewVerbatim}
type:
GGRS26CombattingGerrymanderingRCV.PaperThieleWeights → ℝ → ℕ → ℕ → ℝ

value:
fun (weights : GGRS26CombattingGerrymanderingRCV.PaperThieleWeights) (partyShare : ℝ) (seats seatCount : ℕ) =>
  partyShare * weights.voterScore seatCount + (1 - partyShare) * weights.voterScore (seats - seatCount)
\end{ReviewVerbatim}

\paragraph{Direct retained dependencies}
\begin{itemize}\raggedright
\item \hyperlink{paper-prerequisite-37b91baa3b4ddc9d}{GGRS26\allowbreak{}Combatting\allowbreak{}Gerrymandering\allowbreak{}RCV.\allowbreak{}Paper\allowbreak{}Thiele\allowbreak{}Weights}
\end{itemize}
\hypertarget{governing-declaration-656ba6456b1c16b1}{}
\subsection*{\small\ttfamily\raggedright GGRS26\allowbreak{}Combatting\allowbreak{}Gerrymandering\allowbreak{}RCV.\allowbreak{}paper\_\allowbreak{}pav\_\allowbreak{}min\_\allowbreak{}argmax}
\textbf{Declaration location:} paper-local
\paragraph{Lean declaration body}
\begin{ReviewVerbatim}
type:
ℕ → ℝ → ℕ → Prop

value:
fun (seatCount : ℕ) (partyShare : ℝ) (seats : ℕ) =>
  GGRS26CombattingGerrymanderingRCV.paper_thiele_party_min_argmax
    GGRS26CombattingGerrymanderingRCV.paper_pav_thiele_weights seatCount partyShare seats
\end{ReviewVerbatim}

\paragraph{Direct retained dependencies}
\begin{itemize}\raggedright
\item \hyperlink{governing-declaration-2864a4aebe1e6e50}{GGRS26\allowbreak{}Combatting\allowbreak{}Gerrymandering\allowbreak{}RCV.\allowbreak{}paper\_\allowbreak{}pav\_\allowbreak{}thiele\_\allowbreak{}weights}
\item \hyperlink{paper-prerequisite-58801459d14912d4}{GGRS26\allowbreak{}Combatting\allowbreak{}Gerrymandering\allowbreak{}RCV.\allowbreak{}paper\_\allowbreak{}thiele\_\allowbreak{}party\_\allowbreak{}min\_\allowbreak{}argmax}
\end{itemize}
\hypertarget{governing-declaration-6c3e49fcf5e23bf7}{}
\subsection*{\small\ttfamily\raggedright GGRS26\allowbreak{}Combatting\allowbreak{}Gerrymandering\allowbreak{}RCV.\allowbreak{}pav\allowbreak{}Seat\allowbreak{}Integer\allowbreak{}Interval}
\textbf{Declaration location:} paper-local
\paragraph{Lean declaration body}
\begin{ReviewVerbatim}
type:
ℤ → ℝ → ℕ → Prop

value:
fun (seatCount : ℤ) (partyShare : ℝ) (seats : ℕ) =>
  partyShare * ↑(seats + 1) - 1 ≤ ↑seatCount ∧ ↑seatCount < partyShare * ↑(seats + 1)
\end{ReviewVerbatim}


\hypertarget{governing-declaration-71aa9e7dd85c440d}{}
\subsection*{\small\ttfamily\raggedright GGRS26\allowbreak{}Combatting\allowbreak{}Gerrymandering\allowbreak{}RCV.\allowbreak{}paper\_\allowbreak{}pav\_\allowbreak{}integer\_\allowbreak{}interval}
\textbf{Declaration location:} paper-local
\paragraph{Lean declaration body}
\begin{ReviewVerbatim}
type:
ℤ → ℝ → ℕ → Prop

value:
fun (seatCount : ℤ) (partyShare : ℝ) (seats : ℕ) =>
  GGRS26CombattingGerrymanderingRCV.pavSeatIntegerInterval seatCount partyShare seats
\end{ReviewVerbatim}

\paragraph{Direct retained dependencies}
\begin{itemize}\raggedright
\item \hyperlink{governing-declaration-6c3e49fcf5e23bf7}{GGRS26\allowbreak{}Combatting\allowbreak{}Gerrymandering\allowbreak{}RCV.\allowbreak{}pav\allowbreak{}Seat\allowbreak{}Integer\allowbreak{}Interval}
\end{itemize}
\hypertarget{governing-declaration-d36d27b8173a66a1}{}
\subsection*{\small\ttfamily\raggedright GGRS26\allowbreak{}Combatting\allowbreak{}Gerrymandering\allowbreak{}RCV.\allowbreak{}seat\allowbreak{}Share\allowbreak{}Rounded}
\textbf{Declaration location:} paper-local
\paragraph{Lean declaration body}
\begin{ReviewVerbatim}
type:
ℕ → ℝ → ℕ → Prop

value:
fun (seatCount : ℕ) (partyShare : ℝ) (seats : ℕ) =>
  AppliedModelingLib.SocialChoice.Voting.roundedSeatShare seatCount partyShare seats
\end{ReviewVerbatim}

\paragraph{Direct retained dependencies}
\begin{itemize}\raggedright
\item \hyperlink{governing-declaration-2d4e6f0580886d40}{Applied\allowbreak{}Modeling\allowbreak{}Lib.\allowbreak{}Social\allowbreak{}Choice.\allowbreak{}Voting.\allowbreak{}rounded\allowbreak{}Seat\allowbreak{}Share}
\end{itemize}
\hypertarget{governing-declaration-64adc81e0c40ba7a}{}
\subsection*{\small\ttfamily\raggedright GGRS26\allowbreak{}Combatting\allowbreak{}Gerrymandering\allowbreak{}RCV.\allowbreak{}paper\_\allowbreak{}seat\_\allowbreak{}share\_\allowbreak{}rounded}
\textbf{Declaration location:} paper-local
\paragraph{Lean declaration body}
\begin{ReviewVerbatim}
type:
ℕ → ℝ → ℕ → Prop

value:
fun (seatCount : ℕ) (partyShare : ℝ) (seats : ℕ) =>
  GGRS26CombattingGerrymanderingRCV.seatShareRounded seatCount partyShare seats
\end{ReviewVerbatim}

\paragraph{Direct retained dependencies}
\begin{itemize}\raggedright
\item \hyperlink{governing-declaration-d36d27b8173a66a1}{GGRS26\allowbreak{}Combatting\allowbreak{}Gerrymandering\allowbreak{}RCV.\allowbreak{}seat\allowbreak{}Share\allowbreak{}Rounded}
\end{itemize}
\clearpage
\hypertarget{paper-prerequisite-section-5ef5ef0364b6939c}{}
\section*{Paper-specific semantic prerequisites}
\hypertarget{paper-prerequisite-9acd2440ce1052db}{}
\subsection*{\small\ttfamily\raggedright GGRS26\allowbreak{}Combatting\allowbreak{}Gerrymandering\allowbreak{}RCV.\allowbreak{}Paper\allowbreak{}Source\allowbreak{}STV\allowbreak{}State}
\noindent\textbf{Paper-local declaration:}\par
{\footnotesize\ttfamily\raggedright GGRS26\allowbreak{}Combatting\allowbreak{}Gerrymandering\allowbreak{}RCV.\allowbreak{}Paper\allowbreak{}Source\allowbreak{}STV\allowbreak{}State\par}
\textbf{Source locator:} {\footnotesize\raggedright cited publication, lines 394-\allowbreak{}408\par}
\paragraph{Verbatim paper-source connection}
\begin{ReviewVerbatim}
In Single Transferable Vote (STV), a type of ranked choice voting for multiple winners, each
voter instead submits a ranking over candidates (here, we assume the rankings are complete, not
partial). Consider a district with Nk seats and Vk voters; the set of winners is created iteratively, as
follows. The number of first-place votes for each candidate is counted. Any candidate with a number
of votes at least the “Droop” threshold Q = ⌊ NVk k+1 ⌋ + 1 is selected as a winner, and their surplus
votes (number of votes minus Q) are transferred to each voter’s next preference. If no candidate has
enough votes, then the candidate with the least number of votes (with tie-breaking) is eliminated,
and their votes are transferred. The process continues until Nk candidates have been selected, or
the number of remaining candidates equals the number of remaining seats. STV is commonly used
in practice in multi-winner elections.
    Details can differ in how such votes are transferred, but our results in Section 3 do not depend
on the transfer rule; in our simulations in Section 4 and Appendix B, we use the “fractional” STV
rule, in which each voter keeps a fractional number of votes equal to their share of the winning
candidate’s surplus votes (we describe this formally in Section 4). For consistency, for all rules we
assume that ties are broken in favor of candidates from party D, but randomly within each party.5
\end{ReviewVerbatim}



\paragraph{Lean-elaborated paper signature}
\begin{ReviewVerbatim}
field active:
{Voter : Type u_1} →
  {Candidate : Type u_2} → GGRS26CombattingGerrymanderingRCV.PaperSourceSTVState Voter Candidate → Finset Candidate

field elected:
{Voter : Type u_1} →
  {Candidate : Type u_2} → GGRS26CombattingGerrymanderingRCV.PaperSourceSTVState Voter Candidate → Finset Candidate

field weight:
{Voter : Type u_1} →
  {Candidate : Type u_2} → GGRS26CombattingGerrymanderingRCV.PaperSourceSTVState Voter Candidate → Voter → ℝ
\end{ReviewVerbatim}


\paragraph{Recorded source-to-declaration screening}
\textbf{Verdict:} \texttt{matches}
\\
{\raggedright
\textbf{Reason:} Exposes the active and elected candidate sets and current ballot weights that comprise the source STV count state;\allowbreak{} its execution predicate supplies the associated operational constraints.\allowbreak{}
\par}
\reviewmatch{Matches source input}
\noindent\textbf{Reviewer annotation}\par
\reviewerbox
\clearpage
\hypertarget{paper-prerequisite-37b91baa3b4ddc9d}{}
\subsection*{\small\ttfamily\raggedright GGRS26\allowbreak{}Combatting\allowbreak{}Gerrymandering\allowbreak{}RCV.\allowbreak{}Paper\allowbreak{}Thiele\allowbreak{}Weights}
\noindent\textbf{Paper-local declaration:}\par
{\footnotesize\ttfamily\raggedright GGRS26\allowbreak{}Combatting\allowbreak{}Gerrymandering\allowbreak{}RCV.\allowbreak{}Paper\allowbreak{}Thiele\allowbreak{}Weights\par}
\textbf{Source locator:} {\footnotesize\raggedright cited publication, lines 377-\allowbreak{}393\par}
\paragraph{Verbatim paper-source connection}
\begin{ReviewVerbatim}
A Thiele rule is characterized by a function λ. Each voter v in district k provides a set Sv
of the Nk candidates they like most. The set of winners is determined as follows. Consider a
potential set of winners C. The amount of points that voter v contributes to C is i=1         λ(i), and
                                                                                     P|Sv ∩C|

the set C with the most points across voters (after tie-breaking) is selected. The (non-increasing)
function λ establishes that votes have diminishing returns; the ith approved candidate in the set is
worth λ(i). For example, winner-take-all approval voting is a Thiele rule with λ(i) = 1: for each
voter, the set C receives a number of points equal to the number of candidates in the set that the
voter likes, and so the Nk candidates with the most votes win. We further consider Proportional
Approval Voting (PAV), a Thiele rule with λPAV (i) = 1i ; and Thiele Squared, with λTS = i12 . PAV
is well-studied in the theoretical social choice literature and comes with optimality guarantees (in
terms of proportionality) (Skowron 2021). Thiele Squared induces even sharper diminishing returns
than PAV in a voter having more of their preferred candidates winning, and thus prioritizes more
voters having at least one approved candidate (and 50/50 outcomes between two parties, regardless
of underlying vote shares). More generally, Thiele rules both contain common voting rules and
parameterize a rule’s tendency to favor proportional or balanced outcomes.
\end{ReviewVerbatim}



\paragraph{Lean-elaborated paper signature}
\begin{ReviewVerbatim}
field marginalWeight:
GGRS26CombattingGerrymanderingRCV.PaperThieleWeights → ℕ → ℝ

field marginal_nonincreasing:
∀ (self : GGRS26CombattingGerrymanderingRCV.PaperThieleWeights) (i j : ℕ),
  1 ≤ i → i ≤ j → self.marginalWeight j ≤ self.marginalWeight i

field voterScore:
GGRS26CombattingGerrymanderingRCV.PaperThieleWeights → ℕ → ℝ

field voterScore_zero:
∀ (self : GGRS26CombattingGerrymanderingRCV.PaperThieleWeights), self.voterScore 0 = 0

field voterScore_succ:
∀ (self : GGRS26CombattingGerrymanderingRCV.PaperThieleWeights) (n : ℕ),
  self.voterScore (n + 1) = self.voterScore n + self.marginalWeight (n + 1)
\end{ReviewVerbatim}


\paragraph{Recorded source-to-declaration screening}
\textbf{Verdict:} \texttt{matches}
\\
{\raggedright
\textbf{Reason:} Its positive-\allowbreak{}index nonincreasing marginals, zero cumulative score, and successor recurrence encode the source Thiele sum from i=1 through n, including diminishing returns.\allowbreak{}
\par}
\reviewmatch{Matches source input}
\noindent\textbf{Reviewer annotation}\par
\reviewerbox
\clearpage
\hypertarget{paper-prerequisite-b12522706ad3f624}{}
\subsection*{\small\ttfamily\raggedright GGRS26\allowbreak{}Combatting\allowbreak{}Gerrymandering\allowbreak{}RCV.\allowbreak{}paper\_\allowbreak{}candidate\_\allowbreak{}ranked\_\allowbreak{}before}
\noindent\textbf{Paper-local declaration:}\par
{\footnotesize\ttfamily\raggedright GGRS26\allowbreak{}Combatting\allowbreak{}Gerrymandering\allowbreak{}RCV.\allowbreak{}paper\_\allowbreak{}candidate\_\allowbreak{}ranked\_\allowbreak{}before\par}
\textbf{Source locator:} {\footnotesize\raggedright cited publication, lines 415-\allowbreak{}426\par}
\paragraph{Verbatim paper-source connection}
\begin{ReviewVerbatim}
Voter assumptions. Next, we require assumptions on how voters rank or approve candidates.
For each Thiele rule, we assume that in each district k there are exactly Nk candidates from each
party, and each voter simply approves all the candidates in their party.6 For STV, we assume that
in each district k there are at least Nk candidates from each party.7 For our primarily empirical
analysis, we further assume that parties are solid coalitions (Dummett 1984): that each voter ranks
all candidates of their party over each candidate of the other party, though intra-party rankings
may vary. These assumptions reflect that party membership is a strong predictor of the party of
the candidate for which a voter votes. Such an assumption, using historical party vote shares to
predict vote shares in future elections, is used throughout the redistricting literature for SMDs; our
additional assumption is that this behavior extends to partisan voting in multi-member elections.
We relax this assumption in Section 4 and study intra-party effects, where voters may disagree on
within-party rankings, in Appendix B.
\end{ReviewVerbatim}



\paragraph{Lean-expanded paper semantic target}
\begin{ReviewVerbatim}
type:
{Candidate : Type u_1} → Candidate → Candidate → AppliedModelingLib.SocialChoice.Voting.Ballot Candidate → Prop

value:
fun {Candidate : Type u_1} (higher lower : Candidate)
    (ballot : AppliedModelingLib.SocialChoice.Voting.Ballot Candidate) =>
  ∃ (earlier : AppliedModelingLib.SocialChoice.Voting.Ballot Candidate) (middle : List Candidate) (suffix :
    List Candidate), ballot = earlier ++ higher :: middle ++ lower :: suffix
\end{ReviewVerbatim}

\paragraph{Direct retained dependencies}
\begin{itemize}\raggedright
\item \hyperlink{governing-declaration-4c0320b08e4fb484}{Applied\allowbreak{}Modeling\allowbreak{}Lib.\allowbreak{}Social\allowbreak{}Choice.\allowbreak{}Voting.\allowbreak{}Ballot}
\end{itemize}
\paragraph{Recorded source-to-declaration screening}
\textbf{Verdict:} \texttt{matches}
\\
{\raggedright
\textbf{Reason:} Its list decomposition requires the higher candidate to occur strictly before the lower candidate, which is the exact ranking relation needed for source solid coalitions.\allowbreak{}
\par}
\reviewmatch{Matches source input}
\noindent\textbf{Reviewer annotation}\par
\reviewerbox
\clearpage
\hypertarget{paper-prerequisite-a48f905958044eb5}{}
\subsection*{\small\ttfamily\raggedright GGRS26\allowbreak{}Combatting\allowbreak{}Gerrymandering\allowbreak{}RCV.\allowbreak{}paper\_\allowbreak{}complete\_\allowbreak{}solid\_\allowbreak{}coalition\_\allowbreak{}rankings}
\noindent\textbf{Paper-local declaration:}\par
{\footnotesize\ttfamily\raggedright GGRS26\allowbreak{}Combatting\allowbreak{}Gerrymandering\allowbreak{}RCV.\allowbreak{}paper\_\allowbreak{}complete\_\allowbreak{}solid\_\allowbreak{}coalition\_\allowbreak{}rankings\par}
\textbf{Source locator:} {\footnotesize\raggedright cited publication, lines 415-\allowbreak{}426\par}
\paragraph{Verbatim paper-source connection}
\begin{ReviewVerbatim}
Voter assumptions. Next, we require assumptions on how voters rank or approve candidates.
For each Thiele rule, we assume that in each district k there are exactly Nk candidates from each
party, and each voter simply approves all the candidates in their party.6 For STV, we assume that
in each district k there are at least Nk candidates from each party.7 For our primarily empirical
analysis, we further assume that parties are solid coalitions (Dummett 1984): that each voter ranks
all candidates of their party over each candidate of the other party, though intra-party rankings
may vary. These assumptions reflect that party membership is a strong predictor of the party of
the candidate for which a voter votes. Such an assumption, using historical party vote shares to
predict vote shares in future elections, is used throughout the redistricting literature for SMDs; our
additional assumption is that this behavior extends to partisan voting in multi-member elections.
We relax this assumption in Section 4 and study intra-party effects, where voters may disagree on
within-party rankings, in Appendix B.
\end{ReviewVerbatim}



\paragraph{Lean-expanded paper semantic target}
\begin{ReviewVerbatim}
type:
{Voter : Type u_1} →
  {Candidate : Type u_2} →
    Finset Voter →
      (Voter → AppliedModelingLib.SocialChoice.Voting.Ballot Candidate) →
        Finset Candidate → Finset Candidate → Finset Candidate → Prop

value:
fun {Voter : Type u_1} {Candidate : Type u_2} (voters : Finset Voter)
    (ballots : Voter → AppliedModelingLib.SocialChoice.Voting.Ballot Candidate)
    (partyCandidates otherPartyCandidates allCandidates : Finset Candidate) =>
  Disjoint partyCandidates otherPartyCandidates ∧
    allCandidates = partyCandidates ∪ otherPartyCandidates ∧
      ∀ voter ∈ voters,
        (ballots voter).Valid ∧
          (∀ (candidate : Candidate), candidate ∈ ballots voter ↔ candidate ∈ allCandidates) ∧
            ∀ same ∈ partyCandidates,
              ∀ other ∈ otherPartyCandidates,
                GGRS26CombattingGerrymanderingRCV.paper_candidate_ranked_before same other (ballots voter)
\end{ReviewVerbatim}

\paragraph{Direct retained dependencies}
\begin{itemize}\raggedright
\item \hyperlink{governing-declaration-4c0320b08e4fb484}{Applied\allowbreak{}Modeling\allowbreak{}Lib.\allowbreak{}Social\allowbreak{}Choice.\allowbreak{}Voting.\allowbreak{}Ballot}
\item \hyperlink{governing-declaration-d09ca390012680b9}{Applied\allowbreak{}Modeling\allowbreak{}Lib.\allowbreak{}Social\allowbreak{}Choice.\allowbreak{}Voting.\allowbreak{}Ballot.\allowbreak{}Valid}
\item \hyperlink{paper-prerequisite-b12522706ad3f624}{GGRS26\allowbreak{}Combatting\allowbreak{}Gerrymandering\allowbreak{}RCV.\allowbreak{}paper\_\allowbreak{}candidate\_\allowbreak{}ranked\_\allowbreak{}before}
\end{itemize}
\paragraph{Recorded source-to-declaration screening}
\textbf{Verdict:} \texttt{matches}
\\
{\raggedright
\textbf{Reason:} Uses disjoint complete two-\allowbreak{}party rosters and requires every valid full ballot to rank each own-\allowbreak{}party candidate above every other-\allowbreak{}party candidate, while leaving within-\allowbreak{}party ordering unrestricted as in the source.\allowbreak{}
\par}
\reviewmatch{Matches source input}
\noindent\textbf{Reviewer annotation}\par
\reviewerbox
\clearpage
\hypertarget{paper-prerequisite-e34992194a8b7f0d}{}
\subsection*{\small\ttfamily\raggedright GGRS26\allowbreak{}Combatting\allowbreak{}Gerrymandering\allowbreak{}RCV.\allowbreak{}paper\_\allowbreak{}first\_\allowbreak{}active\_\allowbreak{}candidate}
\noindent\textbf{Paper-local declaration:}\par
{\footnotesize\ttfamily\raggedright GGRS26\allowbreak{}Combatting\allowbreak{}Gerrymandering\allowbreak{}RCV.\allowbreak{}paper\_\allowbreak{}first\_\allowbreak{}active\_\allowbreak{}candidate\par}
\textbf{Source locator:} {\footnotesize\raggedright cited publication, lines 394-\allowbreak{}408\par}
\paragraph{Verbatim paper-source connection}
\begin{ReviewVerbatim}
In Single Transferable Vote (STV), a type of ranked choice voting for multiple winners, each
voter instead submits a ranking over candidates (here, we assume the rankings are complete, not
partial). Consider a district with Nk seats and Vk voters; the set of winners is created iteratively, as
follows. The number of first-place votes for each candidate is counted. Any candidate with a number
of votes at least the “Droop” threshold Q = ⌊ NVk k+1 ⌋ + 1 is selected as a winner, and their surplus
votes (number of votes minus Q) are transferred to each voter’s next preference. If no candidate has
enough votes, then the candidate with the least number of votes (with tie-breaking) is eliminated,
and their votes are transferred. The process continues until Nk candidates have been selected, or
the number of remaining candidates equals the number of remaining seats. STV is commonly used
in practice in multi-winner elections.
    Details can differ in how such votes are transferred, but our results in Section 3 do not depend
on the transfer rule; in our simulations in Section 4 and Appendix B, we use the “fractional” STV
rule, in which each voter keeps a fractional number of votes equal to their share of the winning
candidate’s surplus votes (we describe this formally in Section 4). For consistency, for all rules we
assume that ties are broken in favor of candidates from party D, but randomly within each party.5
\end{ReviewVerbatim}



\paragraph{Lean-expanded paper semantic target}
\begin{ReviewVerbatim}
type:
{Candidate : Type u_1} → AppliedModelingLib.SocialChoice.Voting.Ballot Candidate → Finset Candidate → Candidate → Prop

value:
fun {Candidate : Type u_1} (ballot : AppliedModelingLib.SocialChoice.Voting.Ballot Candidate)
    (active : Finset Candidate) (candidate : Candidate) =>
  candidate ∈ active ∧
    ∃ (preceding : AppliedModelingLib.SocialChoice.Voting.Ballot Candidate) (suffix : List Candidate),
      ballot = preceding ++ candidate :: suffix ∧ ∀ earlier ∈ preceding, earlier ∉ active
\end{ReviewVerbatim}

\paragraph{Direct retained dependencies}
\begin{itemize}\raggedright
\item \hyperlink{governing-declaration-4c0320b08e4fb484}{Applied\allowbreak{}Modeling\allowbreak{}Lib.\allowbreak{}Social\allowbreak{}Choice.\allowbreak{}Voting.\allowbreak{}Ballot}
\end{itemize}
\paragraph{Recorded source-to-declaration screening}
\textbf{Verdict:} \texttt{matches}
\\
{\raggedright
\textbf{Reason:} Requires an active candidate with every earlier listed candidate inactive, exactly the source rule's first currently active preference.\allowbreak{}
\par}
\reviewmatch{Matches source input}
\noindent\textbf{Reviewer annotation}\par
\reviewerbox
\clearpage
\hypertarget{paper-prerequisite-7086a653b0f05ea7}{}
\subsection*{\small\ttfamily\raggedright GGRS26\allowbreak{}Combatting\allowbreak{}Gerrymandering\allowbreak{}RCV.\allowbreak{}paper\_\allowbreak{}thiele\_\allowbreak{}committee\_\allowbreak{}score}
\noindent\textbf{Paper-local declaration:}\par
{\footnotesize\ttfamily\raggedright GGRS26\allowbreak{}Combatting\allowbreak{}Gerrymandering\allowbreak{}RCV.\allowbreak{}paper\_\allowbreak{}thiele\_\allowbreak{}committee\_\allowbreak{}score\par}
\textbf{Source locator:} {\footnotesize\raggedright cited publication, lines 377-\allowbreak{}393\par}
\paragraph{Verbatim paper-source connection}
\begin{ReviewVerbatim}
A Thiele rule is characterized by a function λ. Each voter v in district k provides a set Sv
of the Nk candidates they like most. The set of winners is determined as follows. Consider a
potential set of winners C. The amount of points that voter v contributes to C is i=1         λ(i), and
                                                                                     P|Sv ∩C|

the set C with the most points across voters (after tie-breaking) is selected. The (non-increasing)
function λ establishes that votes have diminishing returns; the ith approved candidate in the set is
worth λ(i). For example, winner-take-all approval voting is a Thiele rule with λ(i) = 1: for each
voter, the set C receives a number of points equal to the number of candidates in the set that the
voter likes, and so the Nk candidates with the most votes win. We further consider Proportional
Approval Voting (PAV), a Thiele rule with λPAV (i) = 1i ; and Thiele Squared, with λTS = i12 . PAV
is well-studied in the theoretical social choice literature and comes with optimality guarantees (in
terms of proportionality) (Skowron 2021). Thiele Squared induces even sharper diminishing returns
than PAV in a voter having more of their preferred candidates winning, and thus prioritizes more
voters having at least one approved candidate (and 50/50 outcomes between two parties, regardless
of underlying vote shares). More generally, Thiele rules both contain common voting rules and
parameterize a rule’s tendency to favor proportional or balanced outcomes.
\end{ReviewVerbatim}



\paragraph{Lean-expanded paper semantic target}
\begin{ReviewVerbatim}
type:
{Candidate : Type u_1} →
  GGRS26CombattingGerrymanderingRCV.PaperThieleWeights →
    Finset Candidate → List (GGRS26CombattingGerrymanderingRCV.PartyApprovalBallot Candidate) → ℝ

value:
fun {Candidate : Type u_1} (weights : GGRS26CombattingGerrymanderingRCV.PaperThieleWeights)
    (committee : Finset Candidate) (profile : List (GGRS26CombattingGerrymanderingRCV.PartyApprovalBallot Candidate)) =>
  AppliedModelingLib.SocialChoice.Voting.thieleScore weights.voterScore committee profile
\end{ReviewVerbatim}

\paragraph{Direct retained dependencies}
\begin{itemize}\raggedright
\item \hyperlink{governing-declaration-6f573ea35df34e74}{Applied\allowbreak{}Modeling\allowbreak{}Lib.\allowbreak{}Social\allowbreak{}Choice.\allowbreak{}Voting.\allowbreak{}thiele\allowbreak{}Score}
\item \hyperlink{paper-prerequisite-37b91baa3b4ddc9d}{GGRS26\allowbreak{}Combatting\allowbreak{}Gerrymandering\allowbreak{}RCV.\allowbreak{}Paper\allowbreak{}Thiele\allowbreak{}Weights}
\item \hyperlink{governing-declaration-3adeede1bc778473}{GGRS26\allowbreak{}Combatting\allowbreak{}Gerrymandering\allowbreak{}RCV.\allowbreak{}Party\allowbreak{}Approval\allowbreak{}Ballot}
\end{itemize}
\paragraph{Recorded source-to-declaration screening}
\textbf{Verdict:} \texttt{matches}
\\
{\raggedright
\textbf{Reason:} Maps each approval set to the cumulative score at its intersection count with the committee and sums across voters, exactly the source Thiele committee-\allowbreak{}score formula.\allowbreak{}
\par}
\reviewmatch{Matches source input}
\noindent\textbf{Reviewer annotation}\par
\reviewerbox
\clearpage
\hypertarget{paper-prerequisite-517db5abc400efba}{}
\subsection*{\small\ttfamily\raggedright GGRS26\allowbreak{}Combatting\allowbreak{}Gerrymandering\allowbreak{}RCV.\allowbreak{}paper\_\allowbreak{}thiele\_\allowbreak{}fixed\_\allowbreak{}size\_\allowbreak{}selected\_\allowbreak{}committee}
\noindent\textbf{Paper-local declaration:}\par
{\footnotesize\ttfamily\raggedright GGRS26\allowbreak{}Combatting\allowbreak{}Gerrymandering\allowbreak{}RCV.\allowbreak{}paper\_\allowbreak{}thiele\_\allowbreak{}fixed\_\allowbreak{}size\_\allowbreak{}selected\_\allowbreak{}committee\par}
\textbf{Source locator:} {\footnotesize\raggedright cited publication, lines 377-\allowbreak{}393\par}
\paragraph{Verbatim paper-source connection}
\begin{ReviewVerbatim}
A Thiele rule is characterized by a function λ. Each voter v in district k provides a set Sv
of the Nk candidates they like most. The set of winners is determined as follows. Consider a
potential set of winners C. The amount of points that voter v contributes to C is i=1         λ(i), and
                                                                                     P|Sv ∩C|

the set C with the most points across voters (after tie-breaking) is selected. The (non-increasing)
function λ establishes that votes have diminishing returns; the ith approved candidate in the set is
worth λ(i). For example, winner-take-all approval voting is a Thiele rule with λ(i) = 1: for each
voter, the set C receives a number of points equal to the number of candidates in the set that the
voter likes, and so the Nk candidates with the most votes win. We further consider Proportional
Approval Voting (PAV), a Thiele rule with λPAV (i) = 1i ; and Thiele Squared, with λTS = i12 . PAV
is well-studied in the theoretical social choice literature and comes with optimality guarantees (in
terms of proportionality) (Skowron 2021). Thiele Squared induces even sharper diminishing returns
than PAV in a voter having more of their preferred candidates winning, and thus prioritizes more
voters having at least one approved candidate (and 50/50 outcomes between two parties, regardless
of underlying vote shares). More generally, Thiele rules both contain common voting rules and
parameterize a rule’s tendency to favor proportional or balanced outcomes.
\end{ReviewVerbatim}



\paragraph{Lean-expanded paper semantic target}
\begin{ReviewVerbatim}
type:
{Candidate : Type u_1} →
  GGRS26CombattingGerrymanderingRCV.PaperThieleWeights →
    (Finset Candidate → Finset Candidate → Prop) →
      Finset Candidate →
        Finset Candidate → List (GGRS26CombattingGerrymanderingRCV.PartyApprovalBallot Candidate) → ℕ → Prop

value:
fun {Candidate : Type u_1} (weights : GGRS26CombattingGerrymanderingRCV.PaperThieleWeights)
    (tiePreferred : Finset Candidate → Finset Candidate → Prop) (candidates committee : Finset Candidate)
    (profile : List (GGRS26CombattingGerrymanderingRCV.PartyApprovalBallot Candidate)) (seats : ℕ) =>
  committee ⊆ candidates ∧
    committee.card = seats ∧
      (∀ other ⊆ candidates,
          other.card = seats →
            GGRS26CombattingGerrymanderingRCV.paper_thiele_committee_score weights other profile ≤
              GGRS26CombattingGerrymanderingRCV.paper_thiele_committee_score weights committee profile) ∧
        ∀ other ⊆ candidates,
          other.card = seats →
            GGRS26CombattingGerrymanderingRCV.paper_thiele_committee_score weights other profile =
                GGRS26CombattingGerrymanderingRCV.paper_thiele_committee_score weights committee profile →
              tiePreferred committee other
\end{ReviewVerbatim}

\paragraph{Direct retained dependencies}
\begin{itemize}\raggedright
\item \hyperlink{paper-prerequisite-37b91baa3b4ddc9d}{GGRS26\allowbreak{}Combatting\allowbreak{}Gerrymandering\allowbreak{}RCV.\allowbreak{}Paper\allowbreak{}Thiele\allowbreak{}Weights}
\item \hyperlink{governing-declaration-3adeede1bc778473}{GGRS26\allowbreak{}Combatting\allowbreak{}Gerrymandering\allowbreak{}RCV.\allowbreak{}Party\allowbreak{}Approval\allowbreak{}Ballot}
\item \hyperlink{paper-prerequisite-7086a653b0f05ea7}{GGRS26\allowbreak{}Combatting\allowbreak{}Gerrymandering\allowbreak{}RCV.\allowbreak{}paper\_\allowbreak{}thiele\_\allowbreak{}committee\_\allowbreak{}score}
\end{itemize}
\paragraph{Recorded source-to-declaration screening}
\textbf{Verdict:} \texttt{matches}
\\
{\raggedright
\textbf{Reason:} Requires a subset of the candidate roster of the fixed district size with globally maximal Thiele score, resolving equal maximizers by the supplied tie preference as the source specifies.\allowbreak{}
\par}
\reviewmatch{Matches source input}
\noindent\textbf{Reviewer annotation}\par
\reviewerbox
\clearpage
\hypertarget{paper-prerequisite-82e254229a66cc0e}{}
\subsection*{\small\ttfamily\raggedright GGRS26\allowbreak{}Combatting\allowbreak{}Gerrymandering\allowbreak{}RCV.\allowbreak{}paper\_\allowbreak{}pav\_\allowbreak{}fixed\_\allowbreak{}size\_\allowbreak{}selected\_\allowbreak{}committee}
\noindent\textbf{Paper-local declaration:}\par
{\footnotesize\ttfamily\raggedright GGRS26\allowbreak{}Combatting\allowbreak{}Gerrymandering\allowbreak{}RCV.\allowbreak{}paper\_\allowbreak{}pav\_\allowbreak{}fixed\_\allowbreak{}size\_\allowbreak{}selected\_\allowbreak{}committee\par}
\textbf{Source locator:} {\footnotesize\raggedright cited publication, lines 377-\allowbreak{}393\par}
\paragraph{Verbatim paper-source connection}
\begin{ReviewVerbatim}
A Thiele rule is characterized by a function λ. Each voter v in district k provides a set Sv
of the Nk candidates they like most. The set of winners is determined as follows. Consider a
potential set of winners C. The amount of points that voter v contributes to C is i=1         λ(i), and
                                                                                     P|Sv ∩C|

the set C with the most points across voters (after tie-breaking) is selected. The (non-increasing)
function λ establishes that votes have diminishing returns; the ith approved candidate in the set is
worth λ(i). For example, winner-take-all approval voting is a Thiele rule with λ(i) = 1: for each
voter, the set C receives a number of points equal to the number of candidates in the set that the
voter likes, and so the Nk candidates with the most votes win. We further consider Proportional
Approval Voting (PAV), a Thiele rule with λPAV (i) = 1i ; and Thiele Squared, with λTS = i12 . PAV
is well-studied in the theoretical social choice literature and comes with optimality guarantees (in
terms of proportionality) (Skowron 2021). Thiele Squared induces even sharper diminishing returns
than PAV in a voter having more of their preferred candidates winning, and thus prioritizes more
voters having at least one approved candidate (and 50/50 outcomes between two parties, regardless
of underlying vote shares). More generally, Thiele rules both contain common voting rules and
parameterize a rule’s tendency to favor proportional or balanced outcomes.
\end{ReviewVerbatim}



\paragraph{Lean-expanded paper semantic target}
\begin{ReviewVerbatim}
type:
{Candidate : Type u_1} →
  (Finset Candidate → Finset Candidate → Prop) →
    Finset Candidate →
      Finset Candidate → List (GGRS26CombattingGerrymanderingRCV.PartyApprovalBallot Candidate) → ℕ → Prop

value:
fun {Candidate : Type u_1} (tiePreferred : Finset Candidate → Finset Candidate → Prop)
    (candidates committee : Finset Candidate)
    (profile : List (GGRS26CombattingGerrymanderingRCV.PartyApprovalBallot Candidate)) (seats : ℕ) =>
  GGRS26CombattingGerrymanderingRCV.paper_thiele_fixed_size_selected_committee
    GGRS26CombattingGerrymanderingRCV.paper_pav_thiele_weights tiePreferred candidates committee profile seats
\end{ReviewVerbatim}

\paragraph{Direct retained dependencies}
\begin{itemize}\raggedright
\item \hyperlink{governing-declaration-3adeede1bc778473}{GGRS26\allowbreak{}Combatting\allowbreak{}Gerrymandering\allowbreak{}RCV.\allowbreak{}Party\allowbreak{}Approval\allowbreak{}Ballot}
\item \hyperlink{governing-declaration-2864a4aebe1e6e50}{GGRS26\allowbreak{}Combatting\allowbreak{}Gerrymandering\allowbreak{}RCV.\allowbreak{}paper\_\allowbreak{}pav\_\allowbreak{}thiele\_\allowbreak{}weights}
\item \hyperlink{paper-prerequisite-517db5abc400efba}{GGRS26\allowbreak{}Combatting\allowbreak{}Gerrymandering\allowbreak{}RCV.\allowbreak{}paper\_\allowbreak{}thiele\_\allowbreak{}fixed\_\allowbreak{}size\_\allowbreak{}selected\_\allowbreak{}committee}
\end{itemize}
\paragraph{Recorded source-to-declaration screening}
\textbf{Verdict:} \texttt{matches}
\\
{\raggedright
\textbf{Reason:} Specializes fixed-\allowbreak{}cardinality global Thiele maximization to harmonic PAV weights and the supplied tie rule, matching the source PAV winner definition.\allowbreak{}
\par}
\reviewmatch{Matches source input}
\noindent\textbf{Reviewer annotation}\par
\reviewerbox
\clearpage
\hypertarget{paper-prerequisite-58801459d14912d4}{}
\subsection*{\small\ttfamily\raggedright GGRS26\allowbreak{}Combatting\allowbreak{}Gerrymandering\allowbreak{}RCV.\allowbreak{}paper\_\allowbreak{}thiele\_\allowbreak{}party\_\allowbreak{}min\_\allowbreak{}argmax}
\noindent\textbf{Paper-local declaration:}\par
{\footnotesize\ttfamily\raggedright GGRS26\allowbreak{}Combatting\allowbreak{}Gerrymandering\allowbreak{}RCV.\allowbreak{}paper\_\allowbreak{}thiele\_\allowbreak{}party\_\allowbreak{}min\_\allowbreak{}argmax\par}
\textbf{Source locator:} {\footnotesize\raggedright cited publication, lines 428-\allowbreak{}456\par}
\paragraph{Verbatim paper-source connection}
\begin{ReviewVerbatim}
Calculating seat share given vote shares. Suppose – for a given district with Nk seats and
Vk voters – we know the fraction yp of voters that belong to each party p ∈ {R, D} (vote share).
How do we calculate the number of winners np from each party for a given social choice function?
   Given our assumptions (including solid coalitions), it is straightforward to efficiently do so for
Thiele rules, as candidates from a given party receive the same votes—and so we do not have to
consider individual candidates, only the number of winners from each party directly. Then, we have
that the party R seat share is:
                                                                             n
                                                         "           "
                                    nR (yR , λ) = min arg max yR
                                                                             X
                                                                                   λ(i)
                                                     n           n
                                                                             i=1
                                                                                   [U+F8F9 private-use glyph][U+F8F9 private-use glyph]
                                                                      k −n
                                                                     NX
                                                      +(1 − yR )             λ(i)[U+F8FB private-use glyph][U+F8FB private-use glyph] .                               (1)
                                                                      i=1

    The inner arg max comes directly from the rule definition – for a given n, suppose a fraction
yR approves n candidates. Then that will yield yR ni=1 λ(i) points, with the remaining fraction
                                                      P

1 − yR approving Nk − n candidates and yielding an additional (1 − yR ) N     i=1 λ(i) points. The
                                                                           P k −n

outer min simply encodes tie-breaking in favor of party D.8 The party D seat share is nD (yR , λ) ≜
Nk − nR (yR , λ).
\end{ReviewVerbatim}



\paragraph{Lean-expanded paper semantic target}
\begin{ReviewVerbatim}
type:
GGRS26CombattingGerrymanderingRCV.PaperThieleWeights → ℕ → ℝ → ℕ → Prop

value:
fun (weights : GGRS26CombattingGerrymanderingRCV.PaperThieleWeights) (seatCount : ℕ) (partyShare : ℝ) (seats : ℕ) =>
  AppliedModelingLib.SocialChoice.Voting.IsMinArgmaxOn
    (GGRS26CombattingGerrymanderingRCV.paper_thiele_party_seat_score weights partyShare seats) seatCount seats
\end{ReviewVerbatim}

\paragraph{Direct retained dependencies}
\begin{itemize}\raggedright
\item \hyperlink{governing-declaration-d2947cd1d156937d}{Applied\allowbreak{}Modeling\allowbreak{}Lib.\allowbreak{}Social\allowbreak{}Choice.\allowbreak{}Voting.\allowbreak{}Is\allowbreak{}Min\allowbreak{}Argmax\allowbreak{}On}
\item \hyperlink{paper-prerequisite-37b91baa3b4ddc9d}{GGRS26\allowbreak{}Combatting\allowbreak{}Gerrymandering\allowbreak{}RCV.\allowbreak{}Paper\allowbreak{}Thiele\allowbreak{}Weights}
\item \hyperlink{governing-declaration-d8487ebdee5a1786}{GGRS26\allowbreak{}Combatting\allowbreak{}Gerrymandering\allowbreak{}RCV.\allowbreak{}paper\_\allowbreak{}thiele\_\allowbreak{}party\_\allowbreak{}seat\_\allowbreak{}score}
\end{itemize}
\paragraph{Recorded source-to-declaration screening}
\textbf{Verdict:} \texttt{matches}
\\
{\raggedright
\textbf{Reason:} Uses the bounded global party-\allowbreak{}score maximum and chooses the smallest co-\allowbreak{}maximizer, exactly the source outer min arg max that favors D in party ties.\allowbreak{}
\par}
\reviewmatch{Matches source input}
\noindent\textbf{Reviewer annotation}\par
\reviewerbox
\clearpage
\hypertarget{paper-prerequisite-edccbf02c4bd6a6e}{}
\subsection*{\small\ttfamily\raggedright GGRS26\allowbreak{}Combatting\allowbreak{}Gerrymandering\allowbreak{}RCV.\allowbreak{}paper\_\allowbreak{}thiele\_\allowbreak{}party\_\allowbreak{}republican\_\allowbreak{}selection}
\noindent\textbf{Paper-local declaration:}\par
{\footnotesize\ttfamily\raggedright GGRS26\allowbreak{}Combatting\allowbreak{}Gerrymandering\allowbreak{}RCV.\allowbreak{}paper\_\allowbreak{}thiele\_\allowbreak{}party\_\allowbreak{}republican\_\allowbreak{}selection\par}
\textbf{Source locator:} {\footnotesize\raggedright cited publication, lines 428-\allowbreak{}456\par}
\paragraph{Verbatim paper-source connection}
\begin{ReviewVerbatim}
Calculating seat share given vote shares. Suppose – for a given district with Nk seats and
Vk voters – we know the fraction yp of voters that belong to each party p ∈ {R, D} (vote share).
How do we calculate the number of winners np from each party for a given social choice function?
   Given our assumptions (including solid coalitions), it is straightforward to efficiently do so for
Thiele rules, as candidates from a given party receive the same votes—and so we do not have to
consider individual candidates, only the number of winners from each party directly. Then, we have
that the party R seat share is:
                                                                             n
                                                         "           "
                                    nR (yR , λ) = min arg max yR
                                                                             X
                                                                                   λ(i)
                                                     n           n
                                                                             i=1
                                                                                   [U+F8F9 private-use glyph][U+F8F9 private-use glyph]
                                                                      k −n
                                                                     NX
                                                      +(1 − yR )             λ(i)[U+F8FB private-use glyph][U+F8FB private-use glyph] .                               (1)
                                                                      i=1

    The inner arg max comes directly from the rule definition – for a given n, suppose a fraction
yR approves n candidates. Then that will yield yR ni=1 λ(i) points, with the remaining fraction
                                                      P

1 − yR approving Nk − n candidates and yielding an additional (1 − yR ) N     i=1 λ(i) points. The
                                                                           P k −n

outer min simply encodes tie-breaking in favor of party D.8 The party D seat share is nD (yR , λ) ≜
Nk − nR (yR , λ).
\end{ReviewVerbatim}



\paragraph{Lean-expanded paper semantic target}
\begin{ReviewVerbatim}
type:
(weights : GGRS26CombattingGerrymanderingRCV.PaperThieleWeights) →
  (partyShare : ℝ) →
    (seats : ℕ) →
      { seatCount : ℕ //
        GGRS26CombattingGerrymanderingRCV.paper_thiele_party_min_argmax weights seatCount partyShare seats }

value:
fun (weights : GGRS26CombattingGerrymanderingRCV.PaperThieleWeights) (partyShare : ℝ) (seats : ℕ) =>
  ⟨Classical.choose
      (GGRS26CombattingGerrymanderingRCV.paper_thiele_party_republican_selection._proof_1 weights partyShare seats),
    GGRS26CombattingGerrymanderingRCV.paper_thiele_party_republican_selection._proof_2 weights partyShare seats⟩
\end{ReviewVerbatim}

\paragraph{Direct retained dependencies}
\begin{itemize}\raggedright
\item \hyperlink{governing-declaration-d2947cd1d156937d}{Applied\allowbreak{}Modeling\allowbreak{}Lib.\allowbreak{}Social\allowbreak{}Choice.\allowbreak{}Voting.\allowbreak{}Is\allowbreak{}Min\allowbreak{}Argmax\allowbreak{}On}
\item \hyperlink{paper-prerequisite-37b91baa3b4ddc9d}{GGRS26\allowbreak{}Combatting\allowbreak{}Gerrymandering\allowbreak{}RCV.\allowbreak{}Paper\allowbreak{}Thiele\allowbreak{}Weights}
\item \hyperlink{paper-prerequisite-58801459d14912d4}{GGRS26\allowbreak{}Combatting\allowbreak{}Gerrymandering\allowbreak{}RCV.\allowbreak{}paper\_\allowbreak{}thiele\_\allowbreak{}party\_\allowbreak{}min\_\allowbreak{}argmax}
\item \hyperlink{governing-declaration-d8487ebdee5a1786}{GGRS26\allowbreak{}Combatting\allowbreak{}Gerrymandering\allowbreak{}RCV.\allowbreak{}paper\_\allowbreak{}thiele\_\allowbreak{}party\_\allowbreak{}seat\_\allowbreak{}score}
\end{itemize}
\paragraph{Recorded source-to-declaration screening}
\textbf{Verdict:} \texttt{matches}
\\
{\raggedright
\textbf{Reason:} The selection subtype carries exactly the bounded maximum and smallest-\allowbreak{}tie conditions of the source n\_\allowbreak{}R definition, without adding implementation semantics.\allowbreak{}
\par}
\reviewmatch{Matches source input}
\noindent\textbf{Reviewer annotation}\par
\reviewerbox
\clearpage
\hypertarget{paper-prerequisite-5fa565502cdaaed9}{}
\subsection*{\small\ttfamily\raggedright GGRS26\allowbreak{}Combatting\allowbreak{}Gerrymandering\allowbreak{}RCV.\allowbreak{}paper\_\allowbreak{}thiele\_\allowbreak{}party\_\allowbreak{}republican\_\allowbreak{}seat\_\allowbreak{}count}
\noindent\textbf{Paper-local declaration:}\par
{\footnotesize\ttfamily\raggedright GGRS26\allowbreak{}Combatting\allowbreak{}Gerrymandering\allowbreak{}RCV.\allowbreak{}paper\_\allowbreak{}thiele\_\allowbreak{}party\_\allowbreak{}republican\_\allowbreak{}seat\_\allowbreak{}count\par}
\textbf{Source locator:} {\footnotesize\raggedright cited publication, lines 428-\allowbreak{}456\par}
\paragraph{Verbatim paper-source connection}
\begin{ReviewVerbatim}
Calculating seat share given vote shares. Suppose – for a given district with Nk seats and
Vk voters – we know the fraction yp of voters that belong to each party p ∈ {R, D} (vote share).
How do we calculate the number of winners np from each party for a given social choice function?
   Given our assumptions (including solid coalitions), it is straightforward to efficiently do so for
Thiele rules, as candidates from a given party receive the same votes—and so we do not have to
consider individual candidates, only the number of winners from each party directly. Then, we have
that the party R seat share is:
                                                                             n
                                                         "           "
                                    nR (yR , λ) = min arg max yR
                                                                             X
                                                                                   λ(i)
                                                     n           n
                                                                             i=1
                                                                                   [U+F8F9 private-use glyph][U+F8F9 private-use glyph]
                                                                      k −n
                                                                     NX
                                                      +(1 − yR )             λ(i)[U+F8FB private-use glyph][U+F8FB private-use glyph] .                               (1)
                                                                      i=1

    The inner arg max comes directly from the rule definition – for a given n, suppose a fraction
yR approves n candidates. Then that will yield yR ni=1 λ(i) points, with the remaining fraction
                                                      P

1 − yR approving Nk − n candidates and yielding an additional (1 − yR ) N     i=1 λ(i) points. The
                                                                           P k −n

outer min simply encodes tie-breaking in favor of party D.8 The party D seat share is nD (yR , λ) ≜
Nk − nR (yR , λ).
\end{ReviewVerbatim}



\paragraph{Lean-expanded paper semantic target}
\begin{ReviewVerbatim}
type:
GGRS26CombattingGerrymanderingRCV.PaperThieleWeights → ℝ → ℕ → ℕ

value:
fun (weights : GGRS26CombattingGerrymanderingRCV.PaperThieleWeights) (partyShare : ℝ) (seats : ℕ) =>
  ↑(GGRS26CombattingGerrymanderingRCV.paper_thiele_party_republican_selection weights partyShare seats)
\end{ReviewVerbatim}

\paragraph{Direct retained dependencies}
\begin{itemize}\raggedright
\item \hyperlink{paper-prerequisite-37b91baa3b4ddc9d}{GGRS26\allowbreak{}Combatting\allowbreak{}Gerrymandering\allowbreak{}RCV.\allowbreak{}Paper\allowbreak{}Thiele\allowbreak{}Weights}
\item \hyperlink{paper-prerequisite-58801459d14912d4}{GGRS26\allowbreak{}Combatting\allowbreak{}Gerrymandering\allowbreak{}RCV.\allowbreak{}paper\_\allowbreak{}thiele\_\allowbreak{}party\_\allowbreak{}min\_\allowbreak{}argmax}
\item \hyperlink{paper-prerequisite-edccbf02c4bd6a6e}{GGRS26\allowbreak{}Combatting\allowbreak{}Gerrymandering\allowbreak{}RCV.\allowbreak{}paper\_\allowbreak{}thiele\_\allowbreak{}party\_\allowbreak{}republican\_\allowbreak{}selection}
\end{itemize}
\paragraph{Recorded source-to-declaration screening}
\textbf{Verdict:} \texttt{matches}
\\
{\raggedright
\textbf{Reason:} Projects the Republican component of the source least-\allowbreak{}maximizer selection while retaining the feasibility, maximality, and tie-\allowbreak{}breaking properties in the paired specification.\allowbreak{}
\par}
\reviewmatch{Matches source input}
\noindent\textbf{Reviewer annotation}\par
\reviewerbox
\clearpage
\hypertarget{paper-prerequisite-c0f468d3ce4de7e0}{}
\subsection*{\small\ttfamily\raggedright GGRS26\allowbreak{}Combatting\allowbreak{}Gerrymandering\allowbreak{}RCV.\allowbreak{}paper\_\allowbreak{}pav\_\allowbreak{}selected\_\allowbreak{}seat\_\allowbreak{}count}
\noindent\textbf{Paper-local declaration:}\par
{\footnotesize\ttfamily\raggedright GGRS26\allowbreak{}Combatting\allowbreak{}Gerrymandering\allowbreak{}RCV.\allowbreak{}paper\_\allowbreak{}pav\_\allowbreak{}selected\_\allowbreak{}seat\_\allowbreak{}count\par}
\textbf{Source locator:} {\footnotesize\raggedright cited publication, lines 428-\allowbreak{}456\par}
\paragraph{Verbatim paper-source connection}
\begin{ReviewVerbatim}
Calculating seat share given vote shares. Suppose – for a given district with Nk seats and
Vk voters – we know the fraction yp of voters that belong to each party p ∈ {R, D} (vote share).
How do we calculate the number of winners np from each party for a given social choice function?
   Given our assumptions (including solid coalitions), it is straightforward to efficiently do so for
Thiele rules, as candidates from a given party receive the same votes—and so we do not have to
consider individual candidates, only the number of winners from each party directly. Then, we have
that the party R seat share is:
                                                                             n
                                                         "           "
                                    nR (yR , λ) = min arg max yR
                                                                             X
                                                                                   λ(i)
                                                     n           n
                                                                             i=1
                                                                                   [U+F8F9 private-use glyph][U+F8F9 private-use glyph]
                                                                      k −n
                                                                     NX
                                                      +(1 − yR )             λ(i)[U+F8FB private-use glyph][U+F8FB private-use glyph] .                               (1)
                                                                      i=1

    The inner arg max comes directly from the rule definition – for a given n, suppose a fraction
yR approves n candidates. Then that will yield yR ni=1 λ(i) points, with the remaining fraction
                                                      P

1 − yR approving Nk − n candidates and yielding an additional (1 − yR ) N     i=1 λ(i) points. The
                                                                           P k −n

outer min simply encodes tie-breaking in favor of party D.8 The party D seat share is nD (yR , λ) ≜
Nk − nR (yR , λ).
\end{ReviewVerbatim}



\paragraph{Lean-expanded paper semantic target}
\begin{ReviewVerbatim}
type:
ℝ → ℕ → ℕ

value:
fun (partyShare : ℝ) (seats : ℕ) =>
  GGRS26CombattingGerrymanderingRCV.paper_thiele_party_republican_seat_count
    GGRS26CombattingGerrymanderingRCV.paper_pav_thiele_weights partyShare seats
\end{ReviewVerbatim}

\paragraph{Direct retained dependencies}
\begin{itemize}\raggedright
\item \hyperlink{governing-declaration-2864a4aebe1e6e50}{GGRS26\allowbreak{}Combatting\allowbreak{}Gerrymandering\allowbreak{}RCV.\allowbreak{}paper\_\allowbreak{}pav\_\allowbreak{}thiele\_\allowbreak{}weights}
\item \hyperlink{paper-prerequisite-5fa565502cdaaed9}{GGRS26\allowbreak{}Combatting\allowbreak{}Gerrymandering\allowbreak{}RCV.\allowbreak{}paper\_\allowbreak{}thiele\_\allowbreak{}party\_\allowbreak{}republican\_\allowbreak{}seat\_\allowbreak{}count}
\end{itemize}
\paragraph{Recorded source-to-declaration screening}
\textbf{Verdict:} \texttt{matches}
\\
{\raggedright
\textbf{Reason:} Projects the constructed PAV least-\allowbreak{}maximizing R-\allowbreak{}seat count, the source n\_\allowbreak{}R selector used in Proposition 1.\allowbreak{}
\par}
\reviewmatch{Matches source input}
\noindent\textbf{Reviewer annotation}\par
\reviewerbox
\clearpage
\hypertarget{paper-prerequisite-652a33f58da6d128}{}
\subsection*{\small\ttfamily\raggedright GGRS26\allowbreak{}Combatting\allowbreak{}Gerrymandering\allowbreak{}RCV.\allowbreak{}paper\_\allowbreak{}pav\_\allowbreak{}selected\_\allowbreak{}seat\_\allowbreak{}count\_\allowbreak{}spec}
\noindent\textbf{Paper-local declaration:}\par
{\footnotesize\ttfamily\raggedright GGRS26\allowbreak{}Combatting\allowbreak{}Gerrymandering\allowbreak{}RCV.\allowbreak{}paper\_\allowbreak{}pav\_\allowbreak{}selected\_\allowbreak{}seat\_\allowbreak{}count\_\allowbreak{}spec\par}
\textbf{Source locator:} {\footnotesize\raggedright cited publication, lines 428-\allowbreak{}456\par}
\paragraph{Verbatim paper-source connection}
\begin{ReviewVerbatim}
Calculating seat share given vote shares. Suppose – for a given district with Nk seats and
Vk voters – we know the fraction yp of voters that belong to each party p ∈ {R, D} (vote share).
How do we calculate the number of winners np from each party for a given social choice function?
   Given our assumptions (including solid coalitions), it is straightforward to efficiently do so for
Thiele rules, as candidates from a given party receive the same votes—and so we do not have to
consider individual candidates, only the number of winners from each party directly. Then, we have
that the party R seat share is:
                                                                             n
                                                         "           "
                                    nR (yR , λ) = min arg max yR
                                                                             X
                                                                                   λ(i)
                                                     n           n
                                                                             i=1
                                                                                   [U+F8F9 private-use glyph][U+F8F9 private-use glyph]
                                                                      k −n
                                                                     NX
                                                      +(1 − yR )             λ(i)[U+F8FB private-use glyph][U+F8FB private-use glyph] .                               (1)
                                                                      i=1

    The inner arg max comes directly from the rule definition – for a given n, suppose a fraction
yR approves n candidates. Then that will yield yR ni=1 λ(i) points, with the remaining fraction
                                                      P

1 − yR approving Nk − n candidates and yielding an additional (1 − yR ) N     i=1 λ(i) points. The
                                                                           P k −n

outer min simply encodes tie-breaking in favor of party D.8 The party D seat share is nD (yR , λ) ≜
Nk − nR (yR , λ).
\end{ReviewVerbatim}



\paragraph{Lean-elaborated paper signature}
\begin{ReviewVerbatim}
∀ (partyShare : ℝ) (seats : ℕ),
  GGRS26CombattingGerrymanderingRCV.paper_pav_min_argmax
    (GGRS26CombattingGerrymanderingRCV.paper_pav_selected_seat_count partyShare seats) partyShare seats
\end{ReviewVerbatim}

\paragraph{Direct retained dependencies}
\begin{itemize}\raggedright
\item \hyperlink{governing-declaration-656ba6456b1c16b1}{GGRS26\allowbreak{}Combatting\allowbreak{}Gerrymandering\allowbreak{}RCV.\allowbreak{}paper\_\allowbreak{}pav\_\allowbreak{}min\_\allowbreak{}argmax}
\item \hyperlink{paper-prerequisite-c0f468d3ce4de7e0}{GGRS26\allowbreak{}Combatting\allowbreak{}Gerrymandering\allowbreak{}RCV.\allowbreak{}paper\_\allowbreak{}pav\_\allowbreak{}selected\_\allowbreak{}seat\_\allowbreak{}count}
\end{itemize}
\paragraph{Recorded source-to-declaration screening}
\textbf{Verdict:} \texttt{matches}
\\
{\raggedright
\textbf{Reason:} States that the constructed PAV R-\allowbreak{}seat count satisfies the full bounded minimum-\allowbreak{}argmax selector relation from the source display.\allowbreak{}
\par}
\reviewmatch{Matches source input}
\noindent\textbf{Reviewer annotation}\par
\reviewerbox
\clearpage
\hypertarget{paper-prerequisite-42c2df760fca2c30}{}
\subsection*{\small\ttfamily\raggedright GGRS26\allowbreak{}Combatting\allowbreak{}Gerrymandering\allowbreak{}RCV.\allowbreak{}paper\_\allowbreak{}source\_\allowbreak{}stv\_\allowbreak{}execution\_\allowbreak{}terminal}
\noindent\textbf{Paper-local declaration:}\par
{\footnotesize\ttfamily\raggedright GGRS26\allowbreak{}Combatting\allowbreak{}Gerrymandering\allowbreak{}RCV.\allowbreak{}paper\_\allowbreak{}source\_\allowbreak{}stv\_\allowbreak{}execution\_\allowbreak{}terminal\par}
\textbf{Source locator:} {\footnotesize\raggedright cited publication, lines 394-\allowbreak{}408\par}
\paragraph{Verbatim paper-source connection}
\begin{ReviewVerbatim}
In Single Transferable Vote (STV), a type of ranked choice voting for multiple winners, each
voter instead submits a ranking over candidates (here, we assume the rankings are complete, not
partial). Consider a district with Nk seats and Vk voters; the set of winners is created iteratively, as
follows. The number of first-place votes for each candidate is counted. Any candidate with a number
of votes at least the “Droop” threshold Q = ⌊ NVk k+1 ⌋ + 1 is selected as a winner, and their surplus
votes (number of votes minus Q) are transferred to each voter’s next preference. If no candidate has
enough votes, then the candidate with the least number of votes (with tie-breaking) is eliminated,
and their votes are transferred. The process continues until Nk candidates have been selected, or
the number of remaining candidates equals the number of remaining seats. STV is commonly used
in practice in multi-winner elections.
    Details can differ in how such votes are transferred, but our results in Section 3 do not depend
on the transfer rule; in our simulations in Section 4 and Appendix B, we use the “fractional” STV
rule, in which each voter keeps a fractional number of votes equal to their share of the winning
candidate’s surplus votes (we describe this formally in Section 4). For consistency, for all rules we
assume that ties are broken in favor of candidates from party D, but randomly within each party.5
\end{ReviewVerbatim}



\paragraph{Lean-expanded paper semantic target}
\begin{ReviewVerbatim}
type:
{Voter : Type u_1} →
  {Candidate : Type u_2} → ℕ → GGRS26CombattingGerrymanderingRCV.PaperSourceSTVState Voter Candidate → Prop

value:
fun {Voter : Type u_1} {Candidate : Type u_2} (seats : ℕ)
    (state : GGRS26CombattingGerrymanderingRCV.PaperSourceSTVState Voter Candidate) =>
  state.elected.card = seats ∨ state.active.card = seats - state.elected.card
\end{ReviewVerbatim}

\paragraph{Direct retained dependencies}
\begin{itemize}\raggedright
\item \hyperlink{paper-prerequisite-9acd2440ce1052db}{GGRS26\allowbreak{}Combatting\allowbreak{}Gerrymandering\allowbreak{}RCV.\allowbreak{}Paper\allowbreak{}Source\allowbreak{}STV\allowbreak{}State}
\end{itemize}
\paragraph{Recorded source-to-declaration screening}
\textbf{Verdict:} \texttt{matches}
\\
{\raggedright
\textbf{Reason:} States precisely the source terminal cases:\allowbreak{} M elected candidates, or the still-\allowbreak{}active candidates fill the seats that remain.\allowbreak{}
\par}
\reviewmatch{Matches source input}
\noindent\textbf{Reviewer annotation}\par
\reviewerbox
\clearpage
\hypertarget{paper-prerequisite-587c1da71549064c}{}
\subsection*{\small\ttfamily\raggedright GGRS26\allowbreak{}Combatting\allowbreak{}Gerrymandering\allowbreak{}RCV.\allowbreak{}paper\_\allowbreak{}source\_\allowbreak{}stv\_\allowbreak{}tally}
\noindent\textbf{Paper-local declaration:}\par
{\footnotesize\ttfamily\raggedright GGRS26\allowbreak{}Combatting\allowbreak{}Gerrymandering\allowbreak{}RCV.\allowbreak{}paper\_\allowbreak{}source\_\allowbreak{}stv\_\allowbreak{}tally\par}
\textbf{Source locator:} {\footnotesize\raggedright cited publication, lines 394-\allowbreak{}408\par}
\paragraph{Verbatim paper-source connection}
\begin{ReviewVerbatim}
In Single Transferable Vote (STV), a type of ranked choice voting for multiple winners, each
voter instead submits a ranking over candidates (here, we assume the rankings are complete, not
partial). Consider a district with Nk seats and Vk voters; the set of winners is created iteratively, as
follows. The number of first-place votes for each candidate is counted. Any candidate with a number
of votes at least the “Droop” threshold Q = ⌊ NVk k+1 ⌋ + 1 is selected as a winner, and their surplus
votes (number of votes minus Q) are transferred to each voter’s next preference. If no candidate has
enough votes, then the candidate with the least number of votes (with tie-breaking) is eliminated,
and their votes are transferred. The process continues until Nk candidates have been selected, or
the number of remaining candidates equals the number of remaining seats. STV is commonly used
in practice in multi-winner elections.
    Details can differ in how such votes are transferred, but our results in Section 3 do not depend
on the transfer rule; in our simulations in Section 4 and Appendix B, we use the “fractional” STV
rule, in which each voter keeps a fractional number of votes equal to their share of the winning
candidate’s surplus votes (we describe this formally in Section 4). For consistency, for all rules we
assume that ties are broken in favor of candidates from party D, but randomly within each party.5
\end{ReviewVerbatim}



\paragraph{Lean-expanded paper semantic target}
\begin{ReviewVerbatim}
type:
{Voter : Type u_1} →
  {Candidate : Type u_2} →
    Finset Voter →
      (Voter → AppliedModelingLib.SocialChoice.Voting.Ballot Candidate) → Finset Candidate → (Voter → ℝ) → Candidate → ℝ

value:
fun {Voter : Type u_1} {Candidate : Type u_2} (allVoters : Finset Voter)
    (ballots : Voter → AppliedModelingLib.SocialChoice.Voting.Ballot Candidate) (active : Finset Candidate)
    (weight : Voter → ℝ) (candidate : Candidate) =>
  ∑ voter ∈ allVoters with
    GGRS26CombattingGerrymanderingRCV.paper_first_active_candidate (ballots voter) active candidate, weight voter
\end{ReviewVerbatim}

\paragraph{Direct retained dependencies}
\begin{itemize}\raggedright
\item \hyperlink{governing-declaration-4c0320b08e4fb484}{Applied\allowbreak{}Modeling\allowbreak{}Lib.\allowbreak{}Social\allowbreak{}Choice.\allowbreak{}Voting.\allowbreak{}Ballot}
\item \hyperlink{paper-prerequisite-e34992194a8b7f0d}{GGRS26\allowbreak{}Combatting\allowbreak{}Gerrymandering\allowbreak{}RCV.\allowbreak{}paper\_\allowbreak{}first\_\allowbreak{}active\_\allowbreak{}candidate}
\end{itemize}
\paragraph{Recorded source-to-declaration screening}
\textbf{Verdict:} \texttt{matches}
\\
{\raggedright
\textbf{Reason:} Sums the current weights of exactly those ballots whose first active listed candidate is the candidate being tallied, matching the source current tally.\allowbreak{}
\par}
\reviewmatch{Matches source input}
\noindent\textbf{Reviewer annotation}\par
\reviewerbox
\clearpage
\hypertarget{paper-prerequisite-01a70a9b3aa00da3}{}
\subsection*{\small\ttfamily\raggedright GGRS26\allowbreak{}Combatting\allowbreak{}Gerrymandering\allowbreak{}RCV.\allowbreak{}paper\_\allowbreak{}source\_\allowbreak{}stv\_\allowbreak{}d\_\allowbreak{}favoring\_\allowbreak{}transition}
\noindent\textbf{Paper-local declaration:}\par
{\footnotesize\ttfamily\raggedright GGRS26\allowbreak{}Combatting\allowbreak{}Gerrymandering\allowbreak{}RCV.\allowbreak{}paper\_\allowbreak{}source\_\allowbreak{}stv\_\allowbreak{}d\_\allowbreak{}favoring\_\allowbreak{}transition\par}
\textbf{Source locator:} {\footnotesize\raggedright cited publication, lines 394-\allowbreak{}408\par}
\paragraph{Verbatim paper-source connection}
\begin{ReviewVerbatim}
In Single Transferable Vote (STV), a type of ranked choice voting for multiple winners, each
voter instead submits a ranking over candidates (here, we assume the rankings are complete, not
partial). Consider a district with Nk seats and Vk voters; the set of winners is created iteratively, as
follows. The number of first-place votes for each candidate is counted. Any candidate with a number
of votes at least the “Droop” threshold Q = ⌊ NVk k+1 ⌋ + 1 is selected as a winner, and their surplus
votes (number of votes minus Q) are transferred to each voter’s next preference. If no candidate has
enough votes, then the candidate with the least number of votes (with tie-breaking) is eliminated,
and their votes are transferred. The process continues until Nk candidates have been selected, or
the number of remaining candidates equals the number of remaining seats. STV is commonly used
in practice in multi-winner elections.
    Details can differ in how such votes are transferred, but our results in Section 3 do not depend
on the transfer rule; in our simulations in Section 4 and Appendix B, we use the “fractional” STV
rule, in which each voter keeps a fractional number of votes equal to their share of the winning
candidate’s surplus votes (we describe this formally in Section 4). For consistency, for all rules we
assume that ties are broken in favor of candidates from party D, but randomly within each party.5
\end{ReviewVerbatim}



\paragraph{Lean-expanded paper semantic target}
\begin{ReviewVerbatim}
type:
{Voter : Type u_1} →
  {Candidate : Type u_2} →
    Finset Voter →
      (Voter → AppliedModelingLib.SocialChoice.Voting.Ballot Candidate) →
        Finset Candidate →
          ℝ →
            ℕ →
              (Finset Candidate → Candidate → (Voter → ℝ) → (Voter → ℝ) → Prop) →
                (Finset Candidate → Candidate → (Voter → ℝ) → (Voter → ℝ) → Prop) →
                  GGRS26CombattingGerrymanderingRCV.PaperSourceSTVState Voter Candidate →
                    GGRS26CombattingGerrymanderingRCV.PaperSourceSTVState Voter Candidate → Prop

value:
fun {Voter : Type u_1} {Candidate : Type u_2} (allVoters : Finset Voter)
    (ballots : Voter → AppliedModelingLib.SocialChoice.Voting.Ballot Candidate) (favoredParty : Finset Candidate)
    (quota : ℝ) (seats : ℕ)
    (electUpdate eliminateUpdate : Finset Candidate → Candidate → (Voter → ℝ) → (Voter → ℝ) → Prop)
    (before after : GGRS26CombattingGerrymanderingRCV.PaperSourceSTVState Voter Candidate) =>
  (∃ (winner : Candidate),
      ¬GGRS26CombattingGerrymanderingRCV.paper_source_stv_execution_terminal seats before ∧
        winner ∈ before.active ∧
          before.elected.card < seats ∧
            quota ≤
                GGRS26CombattingGerrymanderingRCV.paper_source_stv_tally allVoters ballots before.active before.weight
                  winner ∧
              (winner ∉ favoredParty →
                  ∀ candidate ∈ before.active,
                    GGRS26CombattingGerrymanderingRCV.paper_source_stv_tally allVoters ballots before.active
                          before.weight candidate =
                        GGRS26CombattingGerrymanderingRCV.paper_source_stv_tally allVoters ballots before.active
                          before.weight winner →
                      candidate ∉ favoredParty) ∧
                after.active = before.active.erase winner ∧
                  after.elected = insert winner before.elected ∧
                    electUpdate before.active winner before.weight after.weight) ∨
    ∃ (loser : Candidate),
      ¬GGRS26CombattingGerrymanderingRCV.paper_source_stv_execution_terminal seats before ∧
        loser ∈ before.active ∧
          (∀ candidate ∈ before.active,
              GGRS26CombattingGerrymanderingRCV.paper_source_stv_tally allVoters ballots before.active before.weight
                  candidate <
                quota) ∧
            (∀ candidate ∈ before.active,
                GGRS26CombattingGerrymanderingRCV.paper_source_stv_tally allVoters ballots before.active before.weight
                    loser ≤
                  GGRS26CombattingGerrymanderingRCV.paper_source_stv_tally allVoters ballots before.active before.weight
                    candidate) ∧
              (loser ∈ favoredParty →
                  ∀ candidate ∈ before.active,
                    GGRS26CombattingGerrymanderingRCV.paper_source_stv_tally allVoters ballots before.active
                          before.weight candidate =
                        GGRS26CombattingGerrymanderingRCV.paper_source_stv_tally allVoters ballots before.active
                          before.weight loser →
                      candidate ∈ favoredParty) ∧
                after.active = before.active.erase loser ∧
                  after.elected = before.elected ∧ eliminateUpdate before.active loser before.weight after.weight
\end{ReviewVerbatim}

\paragraph{Direct retained dependencies}
\begin{itemize}\raggedright
\item \hyperlink{governing-declaration-4c0320b08e4fb484}{Applied\allowbreak{}Modeling\allowbreak{}Lib.\allowbreak{}Social\allowbreak{}Choice.\allowbreak{}Voting.\allowbreak{}Ballot}
\item \hyperlink{paper-prerequisite-9acd2440ce1052db}{GGRS26\allowbreak{}Combatting\allowbreak{}Gerrymandering\allowbreak{}RCV.\allowbreak{}Paper\allowbreak{}Source\allowbreak{}STV\allowbreak{}State}
\item \hyperlink{paper-prerequisite-42c2df760fca2c30}{GGRS26\allowbreak{}Combatting\allowbreak{}Gerrymandering\allowbreak{}RCV.\allowbreak{}paper\_\allowbreak{}source\_\allowbreak{}stv\_\allowbreak{}execution\_\allowbreak{}terminal}
\item \hyperlink{paper-prerequisite-587c1da71549064c}{GGRS26\allowbreak{}Combatting\allowbreak{}Gerrymandering\allowbreak{}RCV.\allowbreak{}paper\_\allowbreak{}source\_\allowbreak{}stv\_\allowbreak{}tally}
\end{itemize}
\paragraph{Recorded source-to-declaration screening}
\textbf{Verdict:} \texttt{matches}
\\
{\raggedright
\textbf{Reason:} Gives the source-\allowbreak{}shaped quota-\allowbreak{}election-\allowbreak{}before-\allowbreak{}elimination and minimum-\allowbreak{}tally elimination cases, favors D in cross-\allowbreak{}party ties, and represents all within-\allowbreak{}party random outcomes as pathwise support only.\allowbreak{}
\par}
\reviewmatch{Matches source input}
\noindent\textbf{Reviewer annotation}\par
\reviewerbox
\clearpage
\hypertarget{paper-prerequisite-abed5fd2c2d17fba}{}
\subsection*{\small\ttfamily\raggedright GGRS26\allowbreak{}Combatting\allowbreak{}Gerrymandering\allowbreak{}RCV.\allowbreak{}paper\_\allowbreak{}source\_\allowbreak{}surplus\_\allowbreak{}transfer\_\allowbreak{}policy}
\noindent\textbf{Paper-local declaration:}\par
{\footnotesize\ttfamily\raggedright GGRS26\allowbreak{}Combatting\allowbreak{}Gerrymandering\allowbreak{}RCV.\allowbreak{}paper\_\allowbreak{}source\_\allowbreak{}surplus\_\allowbreak{}transfer\_\allowbreak{}policy\par}
\textbf{Source locator:} {\footnotesize\raggedright cited publication, lines 394-\allowbreak{}408\par}
\paragraph{Verbatim paper-source connection}
\begin{ReviewVerbatim}
In Single Transferable Vote (STV), a type of ranked choice voting for multiple winners, each
voter instead submits a ranking over candidates (here, we assume the rankings are complete, not
partial). Consider a district with Nk seats and Vk voters; the set of winners is created iteratively, as
follows. The number of first-place votes for each candidate is counted. Any candidate with a number
of votes at least the “Droop” threshold Q = ⌊ NVk k+1 ⌋ + 1 is selected as a winner, and their surplus
votes (number of votes minus Q) are transferred to each voter’s next preference. If no candidate has
enough votes, then the candidate with the least number of votes (with tie-breaking) is eliminated,
and their votes are transferred. The process continues until Nk candidates have been selected, or
the number of remaining candidates equals the number of remaining seats. STV is commonly used
in practice in multi-winner elections.
    Details can differ in how such votes are transferred, but our results in Section 3 do not depend
on the transfer rule; in our simulations in Section 4 and Appendix B, we use the “fractional” STV
rule, in which each voter keeps a fractional number of votes equal to their share of the winning
candidate’s surplus votes (we describe this formally in Section 4). For consistency, for all rules we
assume that ties are broken in favor of candidates from party D, but randomly within each party.5
\end{ReviewVerbatim}



\paragraph{Lean-expanded paper semantic target}
\begin{ReviewVerbatim}
type:
{Voter : Type u_1} →
  {Candidate : Type u_2} →
    Finset Voter →
      (Voter → AppliedModelingLib.SocialChoice.Voting.Ballot Candidate) →
        ℝ →
          (Finset Candidate → Candidate → (Voter → ℝ) → (Voter → ℝ) → Prop) →
            (Finset Candidate → Candidate → (Voter → ℝ) → (Voter → ℝ) → Prop) → Prop

value:
fun {Voter : Type u_1} {Candidate : Type u_2} (allVoters : Finset Voter)
    (ballots : Voter → AppliedModelingLib.SocialChoice.Voting.Ballot Candidate) (quota : ℝ)
    (electUpdate eliminateUpdate : Finset Candidate → Candidate → (Voter → ℝ) → (Voter → ℝ) → Prop) =>
  (∀ (active : Finset Candidate) (winner : Candidate) (beforeWeight : Voter → ℝ),
      (∀ voter ∈ allVoters, 0 ≤ beforeWeight voter) →
        winner ∈ active →
          quota ≤
              GGRS26CombattingGerrymanderingRCV.paper_source_stv_tally allVoters ballots active beforeWeight winner →
            ∃ (afterWeight : Voter → ℝ), electUpdate active winner beforeWeight afterWeight) ∧
    (∀ (active : Finset Candidate) (loser : Candidate) (beforeWeight : Voter → ℝ),
        (∀ voter ∈ allVoters, 0 ≤ beforeWeight voter) →
          loser ∈ active → ∃ (afterWeight : Voter → ℝ), eliminateUpdate active loser beforeWeight afterWeight) ∧
      (∀ (active : Finset Candidate) (winner : Candidate) (beforeWeight afterWeight : Voter → ℝ),
          (∀ voter ∈ allVoters, 0 ≤ beforeWeight voter) →
            electUpdate active winner beforeWeight afterWeight → ∀ voter ∈ allVoters, 0 ≤ afterWeight voter) ∧
        (∀ (active : Finset Candidate) (winner : Candidate) (beforeWeight afterWeight : Voter → ℝ),
            (∀ voter ∈ allVoters, 0 ≤ beforeWeight voter) →
              electUpdate active winner beforeWeight afterWeight →
                ∀ voter ∈ allVoters,
                  ¬GGRS26CombattingGerrymanderingRCV.paper_first_active_candidate (ballots voter) active winner →
                    afterWeight voter = beforeWeight voter) ∧
          (∀ (active : Finset Candidate) (winner : Candidate) (beforeWeight afterWeight : Voter → ℝ),
              (∀ voter ∈ allVoters, 0 ≤ beforeWeight voter) →
                electUpdate active winner beforeWeight afterWeight →
                  ∑ voter ∈ allVoters with
                      GGRS26CombattingGerrymanderingRCV.paper_first_active_candidate (ballots voter) active winner,
                      afterWeight voter =
                    ∑ voter ∈ allVoters with
                        GGRS26CombattingGerrymanderingRCV.paper_first_active_candidate (ballots voter) active winner,
                        beforeWeight voter -
                      quota) ∧
            ∀ (active : Finset Candidate) (loser : Candidate) (beforeWeight afterWeight : Voter → ℝ),
              (∀ voter ∈ allVoters, 0 ≤ beforeWeight voter) →
                eliminateUpdate active loser beforeWeight afterWeight →
                  ∀ voter ∈ allVoters, afterWeight voter = beforeWeight voter
\end{ReviewVerbatim}

\paragraph{Direct retained dependencies}
\begin{itemize}\raggedright
\item \hyperlink{governing-declaration-4c0320b08e4fb484}{Applied\allowbreak{}Modeling\allowbreak{}Lib.\allowbreak{}Social\allowbreak{}Choice.\allowbreak{}Voting.\allowbreak{}Ballot}
\item \hyperlink{paper-prerequisite-e34992194a8b7f0d}{GGRS26\allowbreak{}Combatting\allowbreak{}Gerrymandering\allowbreak{}RCV.\allowbreak{}paper\_\allowbreak{}first\_\allowbreak{}active\_\allowbreak{}candidate}
\item \hyperlink{paper-prerequisite-587c1da71549064c}{GGRS26\allowbreak{}Combatting\allowbreak{}Gerrymandering\allowbreak{}RCV.\allowbreak{}paper\_\allowbreak{}source\_\allowbreak{}stv\_\allowbreak{}tally}
\end{itemize}
\paragraph{Recorded source-to-declaration screening}
\textbf{Verdict:} \texttt{matches}
\\
{\raggedright
\textbf{Reason:} Requires a quota winner's supporter mass to fall by exactly the quota while off-\allowbreak{}support mass is fixed and leaves eliminations unchanged, matching the source's stated surplus-\allowbreak{}preserving transfer-\allowbreak{}policy scope.\allowbreak{}
\par}
\reviewmatch{Matches source input}
\noindent\textbf{Reviewer annotation}\par
\reviewerbox
\clearpage
\hypertarget{paper-prerequisite-78828233cae137c4}{}
\subsection*{\small\ttfamily\raggedright GGRS26\allowbreak{}Combatting\allowbreak{}Gerrymandering\allowbreak{}RCV.\allowbreak{}paper\_\allowbreak{}source\_\allowbreak{}stv\_\allowbreak{}rule\_\allowbreak{}execution}
\noindent\textbf{Paper-local declaration:}\par
{\footnotesize\ttfamily\raggedright GGRS26\allowbreak{}Combatting\allowbreak{}Gerrymandering\allowbreak{}RCV.\allowbreak{}paper\_\allowbreak{}source\_\allowbreak{}stv\_\allowbreak{}rule\_\allowbreak{}execution\par}
\textbf{Source locator:} {\footnotesize\raggedright cited publication, lines 394-\allowbreak{}408\par}
\paragraph{Verbatim paper-source connection}
\begin{ReviewVerbatim}
In Single Transferable Vote (STV), a type of ranked choice voting for multiple winners, each
voter instead submits a ranking over candidates (here, we assume the rankings are complete, not
partial). Consider a district with Nk seats and Vk voters; the set of winners is created iteratively, as
follows. The number of first-place votes for each candidate is counted. Any candidate with a number
of votes at least the “Droop” threshold Q = ⌊ NVk k+1 ⌋ + 1 is selected as a winner, and their surplus
votes (number of votes minus Q) are transferred to each voter’s next preference. If no candidate has
enough votes, then the candidate with the least number of votes (with tie-breaking) is eliminated,
and their votes are transferred. The process continues until Nk candidates have been selected, or
the number of remaining candidates equals the number of remaining seats. STV is commonly used
in practice in multi-winner elections.
    Details can differ in how such votes are transferred, but our results in Section 3 do not depend
on the transfer rule; in our simulations in Section 4 and Appendix B, we use the “fractional” STV
rule, in which each voter keeps a fractional number of votes equal to their share of the winning
candidate’s surplus votes (we describe this formally in Section 4). For consistency, for all rules we
assume that ties are broken in favor of candidates from party D, but randomly within each party.5
\end{ReviewVerbatim}



\paragraph{Lean-expanded paper semantic target}
\begin{ReviewVerbatim}
type:
{Voter : Type u_1} →
  {Candidate : Type u_2} →
    Finset Voter →
      (Voter → AppliedModelingLib.SocialChoice.Voting.Ballot Candidate) →
        Finset Candidate →
          Finset Candidate →
            ℕ →
              (Finset Candidate → Candidate → (Voter → ℝ) → (Voter → ℝ) → Prop) →
                (Finset Candidate → Candidate → (Voter → ℝ) → (Voter → ℝ) → Prop) →
                  GGRS26CombattingGerrymanderingRCV.PaperSourceSTVState Voter Candidate → Prop

value:
fun {Voter : Type u_1} {Candidate : Type u_2} (allVoters : Finset Voter)
    (ballots : Voter → AppliedModelingLib.SocialChoice.Voting.Ballot Candidate)
    (initialCandidates favoredParty : Finset Candidate) (seats : ℕ)
    (electUpdate eliminateUpdate : Finset Candidate → Candidate → (Voter → ℝ) → (Voter → ℝ) → Prop)
    (terminal : GGRS26CombattingGerrymanderingRCV.PaperSourceSTVState Voter Candidate) =>
  have initial : GGRS26CombattingGerrymanderingRCV.PaperSourceSTVState Voter Candidate :=
    { active := initialCandidates, elected := ∅, weight := fun (x : Voter) => 1 };
  (∀ voter ∈ allVoters,
      (ballots voter).Valid ∧ ∀ (candidate : Candidate), candidate ∈ ballots voter ↔ candidate ∈ initialCandidates) ∧
    favoredParty ⊆ initialCandidates ∧
      GGRS26CombattingGerrymanderingRCV.paper_source_surplus_transfer_policy allVoters ballots
          (↑(AppliedModelingLib.SocialChoice.Voting.STVQuota seats allVoters.card)) electUpdate eliminateUpdate ∧
        Relation.ReflTransGen
            (GGRS26CombattingGerrymanderingRCV.paper_source_stv_d_favoring_transition allVoters ballots favoredParty
              (↑(AppliedModelingLib.SocialChoice.Voting.STVQuota seats allVoters.card)) seats electUpdate
              eliminateUpdate)
            initial terminal ∧
          GGRS26CombattingGerrymanderingRCV.paper_source_stv_execution_terminal seats terminal
\end{ReviewVerbatim}

\paragraph{Direct retained dependencies}
\begin{itemize}\raggedright
\item \hyperlink{governing-declaration-4c0320b08e4fb484}{Applied\allowbreak{}Modeling\allowbreak{}Lib.\allowbreak{}Social\allowbreak{}Choice.\allowbreak{}Voting.\allowbreak{}Ballot}
\item \hyperlink{governing-declaration-d09ca390012680b9}{Applied\allowbreak{}Modeling\allowbreak{}Lib.\allowbreak{}Social\allowbreak{}Choice.\allowbreak{}Voting.\allowbreak{}Ballot.\allowbreak{}Valid}
\item \hyperlink{governing-declaration-d335879a7207e841}{Applied\allowbreak{}Modeling\allowbreak{}Lib.\allowbreak{}Social\allowbreak{}Choice.\allowbreak{}Voting.\allowbreak{}STV\allowbreak{}Quota}
\item \hyperlink{paper-prerequisite-9acd2440ce1052db}{GGRS26\allowbreak{}Combatting\allowbreak{}Gerrymandering\allowbreak{}RCV.\allowbreak{}Paper\allowbreak{}Source\allowbreak{}STV\allowbreak{}State}
\item \hyperlink{paper-prerequisite-01a70a9b3aa00da3}{GGRS26\allowbreak{}Combatting\allowbreak{}Gerrymandering\allowbreak{}RCV.\allowbreak{}paper\_\allowbreak{}source\_\allowbreak{}stv\_\allowbreak{}d\_\allowbreak{}favoring\_\allowbreak{}transition}
\item \hyperlink{paper-prerequisite-42c2df760fca2c30}{GGRS26\allowbreak{}Combatting\allowbreak{}Gerrymandering\allowbreak{}RCV.\allowbreak{}paper\_\allowbreak{}source\_\allowbreak{}stv\_\allowbreak{}execution\_\allowbreak{}terminal}
\item \hyperlink{paper-prerequisite-abed5fd2c2d17fba}{GGRS26\allowbreak{}Combatting\allowbreak{}Gerrymandering\allowbreak{}RCV.\allowbreak{}paper\_\allowbreak{}source\_\allowbreak{}surplus\_\allowbreak{}transfer\_\allowbreak{}policy}
\end{itemize}
\paragraph{Recorded source-to-declaration screening}
\textbf{Verdict:} \texttt{matches}
\\
{\raggedright
\textbf{Reason:} Combines complete valid rankings, exact Droop quota, initial unit weights, a surplus-\allowbreak{}transfer policy, D-\allowbreak{}favoring transitions, and a terminal state;\allowbreak{} this is the selected source execution route and asserts no probability law.\allowbreak{}
\par}
\reviewmatch{Matches source input}
\noindent\textbf{Reviewer annotation}\par
\reviewerbox
\clearpage
\hypertarget{paper-prerequisite-f7db6c9d33a2a094}{}
\subsection*{\small\ttfamily\raggedright GGRS26\allowbreak{}Combatting\allowbreak{}Gerrymandering\allowbreak{}RCV.\allowbreak{}paper\_\allowbreak{}thiele\_\allowbreak{}party\_\allowbreak{}democratic\_\allowbreak{}seat\_\allowbreak{}count}
\noindent\textbf{Paper-local declaration:}\par
{\footnotesize\ttfamily\raggedright GGRS26\allowbreak{}Combatting\allowbreak{}Gerrymandering\allowbreak{}RCV.\allowbreak{}paper\_\allowbreak{}thiele\_\allowbreak{}party\_\allowbreak{}democratic\_\allowbreak{}seat\_\allowbreak{}count\par}
\textbf{Source locator:} {\footnotesize\raggedright cited publication, lines 428-\allowbreak{}456\par}
\paragraph{Verbatim paper-source connection}
\begin{ReviewVerbatim}
Calculating seat share given vote shares. Suppose – for a given district with Nk seats and
Vk voters – we know the fraction yp of voters that belong to each party p ∈ {R, D} (vote share).
How do we calculate the number of winners np from each party for a given social choice function?
   Given our assumptions (including solid coalitions), it is straightforward to efficiently do so for
Thiele rules, as candidates from a given party receive the same votes—and so we do not have to
consider individual candidates, only the number of winners from each party directly. Then, we have
that the party R seat share is:
                                                                             n
                                                         "           "
                                    nR (yR , λ) = min arg max yR
                                                                             X
                                                                                   λ(i)
                                                     n           n
                                                                             i=1
                                                                                   [U+F8F9 private-use glyph][U+F8F9 private-use glyph]
                                                                      k −n
                                                                     NX
                                                      +(1 − yR )             λ(i)[U+F8FB private-use glyph][U+F8FB private-use glyph] .                               (1)
                                                                      i=1

    The inner arg max comes directly from the rule definition – for a given n, suppose a fraction
yR approves n candidates. Then that will yield yR ni=1 λ(i) points, with the remaining fraction
                                                      P

1 − yR approving Nk − n candidates and yielding an additional (1 − yR ) N     i=1 λ(i) points. The
                                                                           P k −n

outer min simply encodes tie-breaking in favor of party D.8 The party D seat share is nD (yR , λ) ≜
Nk − nR (yR , λ).
\end{ReviewVerbatim}



\paragraph{Lean-expanded paper semantic target}
\begin{ReviewVerbatim}
type:
GGRS26CombattingGerrymanderingRCV.PaperThieleWeights → ℝ → ℕ → ℕ

value:
fun (weights : GGRS26CombattingGerrymanderingRCV.PaperThieleWeights) (partyShare : ℝ) (seats : ℕ) =>
  seats - GGRS26CombattingGerrymanderingRCV.paper_thiele_party_republican_seat_count weights partyShare seats
\end{ReviewVerbatim}

\paragraph{Direct retained dependencies}
\begin{itemize}\raggedright
\item \hyperlink{paper-prerequisite-37b91baa3b4ddc9d}{GGRS26\allowbreak{}Combatting\allowbreak{}Gerrymandering\allowbreak{}RCV.\allowbreak{}Paper\allowbreak{}Thiele\allowbreak{}Weights}
\item \hyperlink{paper-prerequisite-5fa565502cdaaed9}{GGRS26\allowbreak{}Combatting\allowbreak{}Gerrymandering\allowbreak{}RCV.\allowbreak{}paper\_\allowbreak{}thiele\_\allowbreak{}party\_\allowbreak{}republican\_\allowbreak{}seat\_\allowbreak{}count}
\end{itemize}
\paragraph{Recorded source-to-declaration screening}
\textbf{Verdict:} \texttt{matches}
\\
{\raggedright
\textbf{Reason:} Defines the Democratic count as M minus the selected Republican count, exactly the source relation n\_\allowbreak{}D = M -\allowbreak{} n\_\allowbreak{}R.\allowbreak{}
\par}
\reviewmatch{Matches source input}
\noindent\textbf{Reviewer annotation}\par
\reviewerbox
\clearpage
\hypertarget{paper-prerequisite-9ae92c4421f50a56}{}
\subsection*{\small\ttfamily\raggedright GGRS26\allowbreak{}Combatting\allowbreak{}Gerrymandering\allowbreak{}RCV.\allowbreak{}paper\_\allowbreak{}thiele\_\allowbreak{}party\_\allowbreak{}republican\_\allowbreak{}seat\_\allowbreak{}count\_\allowbreak{}spec}
\noindent\textbf{Paper-local declaration:}\par
{\footnotesize\ttfamily\raggedright GGRS26\allowbreak{}Combatting\allowbreak{}Gerrymandering\allowbreak{}RCV.\allowbreak{}paper\_\allowbreak{}thiele\_\allowbreak{}party\_\allowbreak{}republican\_\allowbreak{}seat\_\allowbreak{}count\_\allowbreak{}spec\par}
\textbf{Source locator:} {\footnotesize\raggedright cited publication, lines 428-\allowbreak{}456\par}
\paragraph{Verbatim paper-source connection}
\begin{ReviewVerbatim}
Calculating seat share given vote shares. Suppose – for a given district with Nk seats and
Vk voters – we know the fraction yp of voters that belong to each party p ∈ {R, D} (vote share).
How do we calculate the number of winners np from each party for a given social choice function?
   Given our assumptions (including solid coalitions), it is straightforward to efficiently do so for
Thiele rules, as candidates from a given party receive the same votes—and so we do not have to
consider individual candidates, only the number of winners from each party directly. Then, we have
that the party R seat share is:
                                                                             n
                                                         "           "
                                    nR (yR , λ) = min arg max yR
                                                                             X
                                                                                   λ(i)
                                                     n           n
                                                                             i=1
                                                                                   [U+F8F9 private-use glyph][U+F8F9 private-use glyph]
                                                                      k −n
                                                                     NX
                                                      +(1 − yR )             λ(i)[U+F8FB private-use glyph][U+F8FB private-use glyph] .                               (1)
                                                                      i=1

    The inner arg max comes directly from the rule definition – for a given n, suppose a fraction
yR approves n candidates. Then that will yield yR ni=1 λ(i) points, with the remaining fraction
                                                      P

1 − yR approving Nk − n candidates and yielding an additional (1 − yR ) N     i=1 λ(i) points. The
                                                                           P k −n

outer min simply encodes tie-breaking in favor of party D.8 The party D seat share is nD (yR , λ) ≜
Nk − nR (yR , λ).
\end{ReviewVerbatim}



\paragraph{Lean-elaborated paper signature}
\begin{ReviewVerbatim}
∀ (weights : GGRS26CombattingGerrymanderingRCV.PaperThieleWeights) (partyShare : ℝ) (seats : ℕ),
  GGRS26CombattingGerrymanderingRCV.paper_thiele_party_min_argmax weights
    (GGRS26CombattingGerrymanderingRCV.paper_thiele_party_republican_seat_count weights partyShare seats) partyShare
    seats
\end{ReviewVerbatim}

\paragraph{Direct retained dependencies}
\begin{itemize}\raggedright
\item \hyperlink{paper-prerequisite-37b91baa3b4ddc9d}{GGRS26\allowbreak{}Combatting\allowbreak{}Gerrymandering\allowbreak{}RCV.\allowbreak{}Paper\allowbreak{}Thiele\allowbreak{}Weights}
\item \hyperlink{paper-prerequisite-58801459d14912d4}{GGRS26\allowbreak{}Combatting\allowbreak{}Gerrymandering\allowbreak{}RCV.\allowbreak{}paper\_\allowbreak{}thiele\_\allowbreak{}party\_\allowbreak{}min\_\allowbreak{}argmax}
\item \hyperlink{paper-prerequisite-5fa565502cdaaed9}{GGRS26\allowbreak{}Combatting\allowbreak{}Gerrymandering\allowbreak{}RCV.\allowbreak{}paper\_\allowbreak{}thiele\_\allowbreak{}party\_\allowbreak{}republican\_\allowbreak{}seat\_\allowbreak{}count}
\end{itemize}
\paragraph{Recorded source-to-declaration screening}
\textbf{Verdict:} \texttt{matches}
\\
{\raggedright
\textbf{Reason:} States that the selected Republican count obeys the full source bounded least-\allowbreak{}maximizer relation for the two-\allowbreak{}party Thiele score.\allowbreak{}
\par}
\reviewmatch{Matches source input}
\noindent\textbf{Reviewer annotation}\par
\reviewerbox
\clearpage
\hypertarget{paper-prerequisite-1ce27606523bc987}{}
\subsection*{\small\ttfamily\raggedright GGRS26\allowbreak{}Combatting\allowbreak{}Gerrymandering\allowbreak{}RCV.\allowbreak{}paper\_\allowbreak{}thiele\_\allowbreak{}party\_\allowbreak{}selected\_\allowbreak{}seat\_\allowbreak{}counts}
\noindent\textbf{Paper-local declaration:}\par
{\footnotesize\ttfamily\raggedright GGRS26\allowbreak{}Combatting\allowbreak{}Gerrymandering\allowbreak{}RCV.\allowbreak{}paper\_\allowbreak{}thiele\_\allowbreak{}party\_\allowbreak{}selected\_\allowbreak{}seat\_\allowbreak{}counts\par}
\textbf{Source locator:} {\footnotesize\raggedright cited publication, lines 428-\allowbreak{}456\par}
\paragraph{Verbatim paper-source connection}
\begin{ReviewVerbatim}
Calculating seat share given vote shares. Suppose – for a given district with Nk seats and
Vk voters – we know the fraction yp of voters that belong to each party p ∈ {R, D} (vote share).
How do we calculate the number of winners np from each party for a given social choice function?
   Given our assumptions (including solid coalitions), it is straightforward to efficiently do so for
Thiele rules, as candidates from a given party receive the same votes—and so we do not have to
consider individual candidates, only the number of winners from each party directly. Then, we have
that the party R seat share is:
                                                                             n
                                                         "           "
                                    nR (yR , λ) = min arg max yR
                                                                             X
                                                                                   λ(i)
                                                     n           n
                                                                             i=1
                                                                                   [U+F8F9 private-use glyph][U+F8F9 private-use glyph]
                                                                      k −n
                                                                     NX
                                                      +(1 − yR )             λ(i)[U+F8FB private-use glyph][U+F8FB private-use glyph] .                               (1)
                                                                      i=1

    The inner arg max comes directly from the rule definition – for a given n, suppose a fraction
yR approves n candidates. Then that will yield yR ni=1 λ(i) points, with the remaining fraction
                                                      P

1 − yR approving Nk − n candidates and yielding an additional (1 − yR ) N     i=1 λ(i) points. The
                                                                           P k −n

outer min simply encodes tie-breaking in favor of party D.8 The party D seat share is nD (yR , λ) ≜
Nk − nR (yR , λ).
\end{ReviewVerbatim}



\paragraph{Lean-expanded paper semantic target}
\begin{ReviewVerbatim}
type:
GGRS26CombattingGerrymanderingRCV.PaperThieleWeights → ℝ → ℕ → ℕ → ℕ → Prop

value:
fun (weights : GGRS26CombattingGerrymanderingRCV.PaperThieleWeights) (partyShare : ℝ) (seats nR nD : ℕ) =>
  GGRS26CombattingGerrymanderingRCV.paper_thiele_party_min_argmax weights nR partyShare seats ∧ nD = seats - nR
\end{ReviewVerbatim}

\paragraph{Direct retained dependencies}
\begin{itemize}\raggedright
\item \hyperlink{paper-prerequisite-37b91baa3b4ddc9d}{GGRS26\allowbreak{}Combatting\allowbreak{}Gerrymandering\allowbreak{}RCV.\allowbreak{}Paper\allowbreak{}Thiele\allowbreak{}Weights}
\item \hyperlink{paper-prerequisite-58801459d14912d4}{GGRS26\allowbreak{}Combatting\allowbreak{}Gerrymandering\allowbreak{}RCV.\allowbreak{}paper\_\allowbreak{}thiele\_\allowbreak{}party\_\allowbreak{}min\_\allowbreak{}argmax}
\end{itemize}
\paragraph{Recorded source-to-declaration screening}
\textbf{Verdict:} \texttt{matches}
\\
{\raggedright
\textbf{Reason:} Pairs the source least-\allowbreak{}maximizing Republican count with the complementary Democratic count, matching the two reported party seat shares.\allowbreak{}
\par}
\reviewmatch{Matches source input}
\noindent\textbf{Reviewer annotation}\par
\reviewerbox
\clearpage
\hypertarget{paper-prerequisite-52e1d78ea1832965}{}
\subsection*{\small\ttfamily\raggedright GGRS26\allowbreak{}Combatting\allowbreak{}Gerrymandering\allowbreak{}RCV.\allowbreak{}paper\_\allowbreak{}thiele\_\allowbreak{}party\_\allowbreak{}selected\_\allowbreak{}seat\_\allowbreak{}counts\_\allowbreak{}add}
\noindent\textbf{Paper-local declaration:}\par
{\footnotesize\ttfamily\raggedright GGRS26\allowbreak{}Combatting\allowbreak{}Gerrymandering\allowbreak{}RCV.\allowbreak{}paper\_\allowbreak{}thiele\_\allowbreak{}party\_\allowbreak{}selected\_\allowbreak{}seat\_\allowbreak{}counts\_\allowbreak{}add\par}
\textbf{Source locator:} {\footnotesize\raggedright cited publication, lines 428-\allowbreak{}456\par}
\paragraph{Verbatim paper-source connection}
\begin{ReviewVerbatim}
Calculating seat share given vote shares. Suppose – for a given district with Nk seats and
Vk voters – we know the fraction yp of voters that belong to each party p ∈ {R, D} (vote share).
How do we calculate the number of winners np from each party for a given social choice function?
   Given our assumptions (including solid coalitions), it is straightforward to efficiently do so for
Thiele rules, as candidates from a given party receive the same votes—and so we do not have to
consider individual candidates, only the number of winners from each party directly. Then, we have
that the party R seat share is:
                                                                             n
                                                         "           "
                                    nR (yR , λ) = min arg max yR
                                                                             X
                                                                                   λ(i)
                                                     n           n
                                                                             i=1
                                                                                   [U+F8F9 private-use glyph][U+F8F9 private-use glyph]
                                                                      k −n
                                                                     NX
                                                      +(1 − yR )             λ(i)[U+F8FB private-use glyph][U+F8FB private-use glyph] .                               (1)
                                                                      i=1

    The inner arg max comes directly from the rule definition – for a given n, suppose a fraction
yR approves n candidates. Then that will yield yR ni=1 λ(i) points, with the remaining fraction
                                                      P

1 − yR approving Nk − n candidates and yielding an additional (1 − yR ) N     i=1 λ(i) points. The
                                                                           P k −n

outer min simply encodes tie-breaking in favor of party D.8 The party D seat share is nD (yR , λ) ≜
Nk − nR (yR , λ).
\end{ReviewVerbatim}



\paragraph{Lean-elaborated paper signature}
\begin{ReviewVerbatim}
∀ (weights : GGRS26CombattingGerrymanderingRCV.PaperThieleWeights) (partyShare : ℝ) (seats : ℕ),
  GGRS26CombattingGerrymanderingRCV.paper_thiele_party_republican_seat_count weights partyShare seats +
      GGRS26CombattingGerrymanderingRCV.paper_thiele_party_democratic_seat_count weights partyShare seats =
    seats
\end{ReviewVerbatim}

\paragraph{Direct retained dependencies}
\begin{itemize}\raggedright
\item \hyperlink{paper-prerequisite-37b91baa3b4ddc9d}{GGRS26\allowbreak{}Combatting\allowbreak{}Gerrymandering\allowbreak{}RCV.\allowbreak{}Paper\allowbreak{}Thiele\allowbreak{}Weights}
\item \hyperlink{paper-prerequisite-f7db6c9d33a2a094}{GGRS26\allowbreak{}Combatting\allowbreak{}Gerrymandering\allowbreak{}RCV.\allowbreak{}paper\_\allowbreak{}thiele\_\allowbreak{}party\_\allowbreak{}democratic\_\allowbreak{}seat\_\allowbreak{}count}
\item \hyperlink{paper-prerequisite-5fa565502cdaaed9}{GGRS26\allowbreak{}Combatting\allowbreak{}Gerrymandering\allowbreak{}RCV.\allowbreak{}paper\_\allowbreak{}thiele\_\allowbreak{}party\_\allowbreak{}republican\_\allowbreak{}seat\_\allowbreak{}count}
\end{itemize}
\paragraph{Recorded source-to-declaration screening}
\textbf{Verdict:} \texttt{matches}
\\
{\raggedright
\textbf{Reason:} Proves the paired counts sum to the district magnitude through the selected count's bound, expressing the source complement relation without natural-\allowbreak{}number truncation.\allowbreak{}
\par}
\reviewmatch{Matches source input}
\noindent\textbf{Reviewer annotation}\par
\reviewerbox
\clearpage
\hypertarget{paper-prerequisite-2324cfa92955386c}{}
\subsection*{\small\ttfamily\raggedright GGRS26\allowbreak{}Combatting\allowbreak{}Gerrymandering\allowbreak{}RCV.\allowbreak{}paper\_\allowbreak{}thiele\_\allowbreak{}party\_\allowbreak{}selected\_\allowbreak{}seat\_\allowbreak{}counts\_\allowbreak{}spec}
\noindent\textbf{Paper-local declaration:}\par
{\footnotesize\ttfamily\raggedright GGRS26\allowbreak{}Combatting\allowbreak{}Gerrymandering\allowbreak{}RCV.\allowbreak{}paper\_\allowbreak{}thiele\_\allowbreak{}party\_\allowbreak{}selected\_\allowbreak{}seat\_\allowbreak{}counts\_\allowbreak{}spec\par}
\textbf{Source locator:} {\footnotesize\raggedright cited publication, lines 428-\allowbreak{}456\par}
\paragraph{Verbatim paper-source connection}
\begin{ReviewVerbatim}
Calculating seat share given vote shares. Suppose – for a given district with Nk seats and
Vk voters – we know the fraction yp of voters that belong to each party p ∈ {R, D} (vote share).
How do we calculate the number of winners np from each party for a given social choice function?
   Given our assumptions (including solid coalitions), it is straightforward to efficiently do so for
Thiele rules, as candidates from a given party receive the same votes—and so we do not have to
consider individual candidates, only the number of winners from each party directly. Then, we have
that the party R seat share is:
                                                                             n
                                                         "           "
                                    nR (yR , λ) = min arg max yR
                                                                             X
                                                                                   λ(i)
                                                     n           n
                                                                             i=1
                                                                                   [U+F8F9 private-use glyph][U+F8F9 private-use glyph]
                                                                      k −n
                                                                     NX
                                                      +(1 − yR )             λ(i)[U+F8FB private-use glyph][U+F8FB private-use glyph] .                               (1)
                                                                      i=1

    The inner arg max comes directly from the rule definition – for a given n, suppose a fraction
yR approves n candidates. Then that will yield yR ni=1 λ(i) points, with the remaining fraction
                                                      P

1 − yR approving Nk − n candidates and yielding an additional (1 − yR ) N     i=1 λ(i) points. The
                                                                           P k −n

outer min simply encodes tie-breaking in favor of party D.8 The party D seat share is nD (yR , λ) ≜
Nk − nR (yR , λ).
\end{ReviewVerbatim}



\paragraph{Lean-elaborated paper signature}
\begin{ReviewVerbatim}
∀ (weights : GGRS26CombattingGerrymanderingRCV.PaperThieleWeights) (partyShare : ℝ) (seats : ℕ),
  GGRS26CombattingGerrymanderingRCV.paper_thiele_party_selected_seat_counts weights partyShare seats
    (GGRS26CombattingGerrymanderingRCV.paper_thiele_party_republican_seat_count weights partyShare seats)
    (GGRS26CombattingGerrymanderingRCV.paper_thiele_party_democratic_seat_count weights partyShare seats)
\end{ReviewVerbatim}

\paragraph{Direct retained dependencies}
\begin{itemize}\raggedright
\item \hyperlink{paper-prerequisite-37b91baa3b4ddc9d}{GGRS26\allowbreak{}Combatting\allowbreak{}Gerrymandering\allowbreak{}RCV.\allowbreak{}Paper\allowbreak{}Thiele\allowbreak{}Weights}
\item \hyperlink{paper-prerequisite-f7db6c9d33a2a094}{GGRS26\allowbreak{}Combatting\allowbreak{}Gerrymandering\allowbreak{}RCV.\allowbreak{}paper\_\allowbreak{}thiele\_\allowbreak{}party\_\allowbreak{}democratic\_\allowbreak{}seat\_\allowbreak{}count}
\item \hyperlink{paper-prerequisite-5fa565502cdaaed9}{GGRS26\allowbreak{}Combatting\allowbreak{}Gerrymandering\allowbreak{}RCV.\allowbreak{}paper\_\allowbreak{}thiele\_\allowbreak{}party\_\allowbreak{}republican\_\allowbreak{}seat\_\allowbreak{}count}
\item \hyperlink{paper-prerequisite-1ce27606523bc987}{GGRS26\allowbreak{}Combatting\allowbreak{}Gerrymandering\allowbreak{}RCV.\allowbreak{}paper\_\allowbreak{}thiele\_\allowbreak{}party\_\allowbreak{}selected\_\allowbreak{}seat\_\allowbreak{}counts}
\end{itemize}
\paragraph{Recorded source-to-declaration screening}
\textbf{Verdict:} \texttt{matches}
\\
{\raggedright
\textbf{Reason:} States the exact source selector relation together with the Democratic complement for the constructed pair of party seat counts.\allowbreak{}
\par}
\reviewmatch{Matches source input}
\noindent\textbf{Reviewer annotation}\par
\reviewerbox
\clearpage
\hypertarget{paper-prerequisite-aea2ad46d27b40e1}{}
\subsection*{\small\ttfamily\raggedright GGRS26\allowbreak{}Combatting\allowbreak{}Gerrymandering\allowbreak{}RCV.\allowbreak{}paper\_\allowbreak{}two\_\allowbreak{}party\_\allowbreak{}voter\_\allowbreak{}model}
\noindent\textbf{Paper-local declaration:}\par
{\footnotesize\ttfamily\raggedright GGRS26\allowbreak{}Combatting\allowbreak{}Gerrymandering\allowbreak{}RCV.\allowbreak{}paper\_\allowbreak{}two\_\allowbreak{}party\_\allowbreak{}voter\_\allowbreak{}model\par}
\textbf{Source locator:} {\footnotesize\raggedright cited publication, lines 415-\allowbreak{}426\par}
\paragraph{Verbatim paper-source connection}
\begin{ReviewVerbatim}
Voter assumptions. Next, we require assumptions on how voters rank or approve candidates.
For each Thiele rule, we assume that in each district k there are exactly Nk candidates from each
party, and each voter simply approves all the candidates in their party.6 For STV, we assume that
in each district k there are at least Nk candidates from each party.7 For our primarily empirical
analysis, we further assume that parties are solid coalitions (Dummett 1984): that each voter ranks
all candidates of their party over each candidate of the other party, though intra-party rankings
may vary. These assumptions reflect that party membership is a strong predictor of the party of
the candidate for which a voter votes. Such an assumption, using historical party vote shares to
predict vote shares in future elections, is used throughout the redistricting literature for SMDs; our
additional assumption is that this behavior extends to partisan voting in multi-member elections.
We relax this assumption in Section 4 and study intra-party effects, where voters may disagree on
within-party rankings, in Appendix B.
\end{ReviewVerbatim}



\paragraph{Lean-expanded paper semantic target}
\begin{ReviewVerbatim}
type:
{Voter : Type u_1} →
  {Candidate : Type u_2} →
    Finset Voter →
      Finset Voter →
        Finset Voter →
          Finset Candidate →
            Finset Candidate →
              (Voter → GGRS26CombattingGerrymanderingRCV.PartyApprovalBallot Candidate) →
                Finset Candidate →
                  Finset Candidate →
                    Finset Candidate → (Voter → AppliedModelingLib.SocialChoice.Voting.Ballot Candidate) → ℕ → Prop

value:
fun {Voter : Type u_1} {Candidate : Type u_2} (allVoters partyVoters otherPartyVoters : Finset Voter)
    (thielePartyCandidates thieleOtherPartyCandidates : Finset Candidate)
    (approvalBallots : Voter → GGRS26CombattingGerrymanderingRCV.PartyApprovalBallot Candidate)
    (stvPartyCandidates stvOtherPartyCandidates allSTVCandidates : Finset Candidate)
    (rankingBallots : Voter → AppliedModelingLib.SocialChoice.Voting.Ballot Candidate) (seats : ℕ) =>
  allVoters = partyVoters ∪ otherPartyVoters ∧
    Disjoint partyVoters otherPartyVoters ∧
      Disjoint thielePartyCandidates thieleOtherPartyCandidates ∧
        thielePartyCandidates.card = seats ∧
          thieleOtherPartyCandidates.card = seats ∧
            (∀ voter ∈ partyVoters, approvalBallots voter = thielePartyCandidates) ∧
              (∀ voter ∈ otherPartyVoters, approvalBallots voter = thieleOtherPartyCandidates) ∧
                seats ≤ stvPartyCandidates.card ∧
                  seats ≤ stvOtherPartyCandidates.card ∧
                    GGRS26CombattingGerrymanderingRCV.paper_complete_solid_coalition_rankings partyVoters rankingBallots
                        stvPartyCandidates stvOtherPartyCandidates allSTVCandidates ∧
                      GGRS26CombattingGerrymanderingRCV.paper_complete_solid_coalition_rankings otherPartyVoters
                        rankingBallots stvOtherPartyCandidates stvPartyCandidates allSTVCandidates
\end{ReviewVerbatim}

\paragraph{Direct retained dependencies}
\begin{itemize}\raggedright
\item \hyperlink{governing-declaration-4c0320b08e4fb484}{Applied\allowbreak{}Modeling\allowbreak{}Lib.\allowbreak{}Social\allowbreak{}Choice.\allowbreak{}Voting.\allowbreak{}Ballot}
\item \hyperlink{governing-declaration-3adeede1bc778473}{GGRS26\allowbreak{}Combatting\allowbreak{}Gerrymandering\allowbreak{}RCV.\allowbreak{}Party\allowbreak{}Approval\allowbreak{}Ballot}
\item \hyperlink{paper-prerequisite-a48f905958044eb5}{GGRS26\allowbreak{}Combatting\allowbreak{}Gerrymandering\allowbreak{}RCV.\allowbreak{}paper\_\allowbreak{}complete\_\allowbreak{}solid\_\allowbreak{}coalition\_\allowbreak{}rankings}
\end{itemize}
\paragraph{Recorded source-to-declaration screening}
\textbf{Verdict:} \texttt{matches}
\\
{\raggedright
\textbf{Reason:} Uses disjoint R/\allowbreak{}D voter blocks and candidate rosters, party approval sets, enough candidates for each rule, and complete solid-\allowbreak{}coalition rankings, matching the source two-\allowbreak{}party electorate and ballot model.\allowbreak{}
\par}
\reviewmatch{Matches source input}
\noindent\textbf{Reviewer annotation}\par
\reviewerbox

\clearpage
\hypertarget{source-claim-6ccc405fd526f01f}{}
\hypertarget{source-section-f10b5f68e4acf3dc}{}
\section*{Source claims}
\subsection*{Review row 1}
\textbf{Source locator:} {\footnotesize\raggedright cited publication, lines 2414-\allowbreak{}2433\par}
\paragraph{Verbatim source input}
\begin{ReviewVerbatim}
Lemma C.1. Let yR ∈ (0, 1]. Consider
                                                                          n
                                                       "           "
                              nR (yR , λP AV ) = min arg max yR
                                                                          X
                                                                                λ(i)
                                                  n         n
                                                                          i=1
                                                                   M −n
                                                                                 ##
                                                      +(1 − yR )
                                                                   X
                                                                          λ(i)        .
                                                                   i=1

    and let ell be the unique integer such that

                                    yR (M + 1) − 1 ≤ ell < yR (M + 1)

    Then, nR (yR , λPAV ) = ell, for all yR , M .

[Next verbatim source excerpt]

    A Thiele rule is characterized by a function λ. Each voter v in district k provides a set Sv
of the Nk candidates they like most. The set of winners is determined as follows. Consider a
potential set of winners C. The amount of points that voter v contributes to C is i=1         λ(i), and
                                                                                     P|Sv ∩C|

the set C with the most points across voters (after tie-breaking) is selected. The (non-increasing)
function λ establishes that votes have diminishing returns; the ith approved candidate in the set is
worth λ(i). For example, winner-take-all approval voting is a Thiele rule with λ(i) = 1: for each
voter, the set C receives a number of points equal to the number of candidates in the set that the
voter likes, and so the Nk candidates with the most votes win. We further consider Proportional
Approval Voting (PAV), a Thiele rule with λPAV (i) = 1i ; and Thiele Squared, with λTS = i12 . PAV
is well-studied in the theoretical social choice literature and comes with optimality guarantees (in
terms of proportionality) (Skowron 2021). Thiele Squared induces even sharper diminishing returns
than PAV in a voter having more of their preferred candidates winning, and thus prioritizes more
voters having at least one approved candidate (and 50/50 outcomes between two parties, regardless
of underlying vote shares). More generally, Thiele rules both contain common voting rules and
parameterize a rule’s tendency to favor proportional or balanced outcomes.
\end{ReviewVerbatim}



\paragraph{Expanded PaperInterface specification}
\begin{ReviewVerbatim}
∀ {seats : ℕ} {partyShare : ℝ},
  1 ≤ seats →
    0 < partyShare →
      partyShare ≤ 1 →
        ∃ (seatCount : ℕ),
          GGRS26CombattingGerrymanderingRCV.paper_pav_min_argmax seatCount partyShare seats ∧
            ∃ (ell : ℤ),
              ↑seatCount = ell ∧
                GGRS26CombattingGerrymanderingRCV.paper_pav_integer_interval ell partyShare seats ∧
                  (∀ (other : ℤ),
                      GGRS26CombattingGerrymanderingRCV.paper_pav_integer_interval other partyShare seats →
                        other = ell) ∧
                    ∀ (otherSeatCount : ℕ),
                      GGRS26CombattingGerrymanderingRCV.paper_pav_min_argmax otherSeatCount partyShare seats →
                        ↑otherSeatCount = ell
\end{ReviewVerbatim}

\paragraph{Retained governing declarations}
\begin{itemize}\raggedright
\item \hyperlink{governing-declaration-71aa9e7dd85c440d}{GGRS26\allowbreak{}Combatting\allowbreak{}Gerrymandering\allowbreak{}RCV.\allowbreak{}paper\_\allowbreak{}pav\_\allowbreak{}integer\_\allowbreak{}interval}
\item \hyperlink{governing-declaration-656ba6456b1c16b1}{GGRS26\allowbreak{}Combatting\allowbreak{}Gerrymandering\allowbreak{}RCV.\allowbreak{}paper\_\allowbreak{}pav\_\allowbreak{}min\_\allowbreak{}argmax}
\end{itemize}
\paragraph{Lean-elaborated claim roles}\begin{ReviewVerbatim}
1. parameter: ℕ
2. parameter: ℝ
3. assumption: 1 ≤ seats
4. assumption: 0 < partyShare
5. assumption: partyShare ≤ 1
6. conclusion: ∃ (seatCount : ℕ),
  GGRS26CombattingGerrymanderingRCV.paper_pav_min_argmax seatCount partyShare seats ∧
    ∃ (ell : ℤ),
      ↑seatCount = ell ∧
        GGRS26CombattingGerrymanderingRCV.paper_pav_integer_interval ell partyShare seats ∧
          (∀ (other : ℤ),
              GGRS26CombattingGerrymanderingRCV.paper_pav_integer_interval other partyShare seats → other = ell) ∧
            ∀ (otherSeatCount : ℕ),
              GGRS26CombattingGerrymanderingRCV.paper_pav_min_argmax otherSeatCount partyShare seats →
                ↑otherSeatCount = ell
\end{ReviewVerbatim}

\paragraph{Lean proof endpoint (built separately)}\begin{ReviewVerbatim}
GGRS26CombattingGerrymanderingRCV.paper_lemma_c1_pav_selector_eq_unique_integer_interval
\end{ReviewVerbatim}

\paragraph{Recorded source-to-Spec screening}
\textbf{Verdict:} \texttt{matches}
\\
{\raggedright
\textbf{Reason:} The frozen target has the source's positive party-\allowbreak{}share and positive district-\allowbreak{}magnitude domain, PAV harmonic leftmost maximizer on 0 through M, the displayed half-\allowbreak{}open integer interval, and equality of every such selector to its unique ell.\allowbreak{}
\par}
\reviewmatch{Matches source input}
\noindent\textbf{Reviewer annotation}\par
\reviewerbox
\clearpage
\hypertarget{source-claim-1d81d5a90a327c0b}{}
\section*{Review row 2}
\textbf{Source locator:} {\footnotesize\raggedright cited publication, lines 485-\allowbreak{}490\par}
\paragraph{Verbatim source input}
\begin{ReviewVerbatim}
Proposition 1 (Seat shares under STV). Suppose – for a given district with M seats and V ≥
M (M + 1) voters – that a fraction yp of voters belong to each party p ∈ {R, D}, and that there are
at least M candidates per party. Assume that each party’s voters rank all same-party candidates
above all other-party candidates, and ties are broken in party D’s favor.
    Let nR (yR , ST V ) be the number of winners belonging to party R using STV. Then, both
nR (yR , ST V ) and nR (yR , λPAV ) are in {⌊yR M ⌋, ⌈yR M ⌉}.

[Next verbatim source excerpt]

    In Single Transferable Vote (STV), a type of ranked choice voting for multiple winners, each
voter instead submits a ranking over candidates (here, we assume the rankings are complete, not
partial). Consider a district with Nk seats and Vk voters; the set of winners is created iteratively, as
follows. The number of first-place votes for each candidate is counted. Any candidate with a number
of votes at least the “Droop” threshold Q = ⌊ NVk k+1 ⌋ + 1 is selected as a winner, and their surplus
votes (number of votes minus Q) are transferred to each voter’s next preference. If no candidate has
enough votes, then the candidate with the least number of votes (with tie-breaking) is eliminated,
and their votes are transferred. The process continues until Nk candidates have been selected, or
the number of remaining candidates equals the number of remaining seats. STV is commonly used
in practice in multi-winner elections.
    Details can differ in how such votes are transferred, but our results in Section 3 do not depend
on the transfer rule; in our simulations in Section 4 and Appendix B, we use the “fractional” STV
rule, in which each voter keeps a fractional number of votes equal to their share of the winning
candidate’s surplus votes (we describe this formally in Section 4). For consistency, for all rules we
assume that ties are broken in favor of candidates from party D, but randomly within each party.5

[Next verbatim source excerpt]

    At a high level, the proof establishes that no candidate receives meaningful (in terms of party
vote share) votes from voters of the other party; in terms of party vote share, we can proceed with
running two separate STV processes. Then, the sequential round path dependence does not matter,
as the overall number of first-place votes for candidates of each party is invariant, and so the final
partisan outcome depends just on initial vote shares. We note that the proof does not depend on
our assumption of fractional STV; other transfer rules (such as random voter transfer) that preserve
the number of surplus votes would induce seat shares according to the same formula.

[Next verbatim source excerpt]

    A Thiele rule is characterized by a function λ. Each voter v in district k provides a set Sv
of the Nk candidates they like most. The set of winners is determined as follows. Consider a
potential set of winners C. The amount of points that voter v contributes to C is i=1         λ(i), and
                                                                                     P|Sv ∩C|

the set C with the most points across voters (after tie-breaking) is selected. The (non-increasing)
function λ establishes that votes have diminishing returns; the ith approved candidate in the set is
worth λ(i). For example, winner-take-all approval voting is a Thiele rule with λ(i) = 1: for each
voter, the set C receives a number of points equal to the number of candidates in the set that the
voter likes, and so the Nk candidates with the most votes win. We further consider Proportional
Approval Voting (PAV), a Thiele rule with λPAV (i) = 1i ; and Thiele Squared, with λTS = i12 . PAV
is well-studied in the theoretical social choice literature and comes with optimality guarantees (in
terms of proportionality) (Skowron 2021). Thiele Squared induces even sharper diminishing returns
than PAV in a voter having more of their preferred candidates winning, and thus prioritizes more
voters having at least one approved candidate (and 50/50 outcomes between two parties, regardless
of underlying vote shares). More generally, Thiele rules both contain common voting rules and
parameterize a rule’s tendency to favor proportional or balanced outcomes.

[Next verbatim source excerpt]

    In Single Transferable Vote (STV), a type of ranked choice voting for multiple winners, each
voter instead submits a ranking over candidates (here, we assume the rankings are complete, not
partial). Consider a district with Nk seats and Vk voters; the set of winners is created iteratively, as
follows. The number of first-place votes for each candidate is counted. Any candidate with a number
of votes at least the “Droop” threshold Q = ⌊ NVk k+1 ⌋ + 1 is selected as a winner, and their surplus
votes (number of votes minus Q) are transferred to each voter’s next preference. If no candidate has
enough votes, then the candidate with the least number of votes (with tie-breaking) is eliminated,
and their votes are transferred. The process continues until Nk candidates have been selected, or
the number of remaining candidates equals the number of remaining seats. STV is commonly used
in practice in multi-winner elections.
    Details can differ in how such votes are transferred, but our results in Section 3 do not depend
on the transfer rule; in our simulations in Section 4 and Appendix B, we use the “fractional” STV
rule, in which each voter keeps a fractional number of votes equal to their share of the winning
candidate’s surplus votes (we describe this formally in Section 4). For consistency, for all rules we
assume that ties are broken in favor of candidates from party D, but randomly within each party.5
\end{ReviewVerbatim}



\paragraph{Expanded PaperInterface specification}
\begin{ReviewVerbatim}
∀ {Voter : Type u_1} {Candidate : Type u_2} {partyVoters otherPartyVoters allVoters : Finset Voter}
  {ballots : Voter → AppliedModelingLib.SocialChoice.Voting.Ballot Candidate}
  {partyCandidates otherPartyCandidates initialActive : Finset Candidate} {seats voters : ℕ} {partyShare : ℝ},
  1 ≤ seats →
    0 ≤ partyShare →
      partyShare ≤ 1 →
        seats * (seats + 1) ≤ voters →
          seats ≤ partyCandidates.card →
            seats ≤ otherPartyCandidates.card →
              GGRS26CombattingGerrymanderingRCV.paper_complete_party_ranking_ballots partyVoters ballots partyCandidates
                  initialActive →
                GGRS26CombattingGerrymanderingRCV.paper_complete_party_ranking_ballots otherPartyVoters ballots
                    otherPartyCandidates initialActive →
                  voters = allVoters.card →
                    allVoters = partyVoters ∪ otherPartyVoters →
                      Disjoint partyVoters otherPartyVoters →
                        Disjoint partyCandidates otherPartyCandidates →
                          partyCandidates ⊆ initialActive →
                            otherPartyCandidates ⊆ initialActive →
                              initialActive ⊆ partyCandidates ∪ otherPartyCandidates →
                                partyShare * ↑voters = ↑partyVoters.card →
                                  (1 - partyShare) * ↑voters = ↑otherPartyVoters.card →
                                    (∀
                                        (policy :
                                          GGRS26CombattingGerrymanderingRCV.BallotRoutedSTVTransferPolicy allVoters
                                            ballots ↑(AppliedModelingLib.SocialChoice.Voting.STVQuota seats voters)),
                                        (∃ (terminal :
                                            GGRS26CombattingGerrymanderingRCV.BallotRoutedSTVState allVoters
                                              initialActive),
                                            GGRS26CombattingGerrymanderingRCV.BallotRoutedDFavoringSTVRun ballots
                                                (↑(AppliedModelingLib.SocialChoice.Voting.STVQuota seats voters)) policy
                                                seats
                                                (GGRS26CombattingGerrymanderingRCV.paper_proposition1_ballot_routed_stv_initial_state
                                                  allVoters initialActive)
                                                terminal ∧
                                              GGRS26CombattingGerrymanderingRCV.BallotRoutedSTVTerminal seats
                                                terminal) ∧
                                          ∀
                                            (terminal :
                                              GGRS26CombattingGerrymanderingRCV.BallotRoutedSTVState allVoters
                                                initialActive),
                                            GGRS26CombattingGerrymanderingRCV.BallotRoutedDFavoringSTVRun ballots
                                                (↑(AppliedModelingLib.SocialChoice.Voting.STVQuota seats voters)) policy
                                                seats
                                                (GGRS26CombattingGerrymanderingRCV.paper_proposition1_ballot_routed_stv_initial_state
                                                  allVoters initialActive)
                                                terminal →
                                              GGRS26CombattingGerrymanderingRCV.BallotRoutedSTVTerminal seats terminal →
                                                GGRS26CombattingGerrymanderingRCV.paper_seat_share_rounded
                                                  (GGRS26CombattingGerrymanderingRCV.ballotRoutedPartyFinalSeats
                                                    partyCandidates seats terminal)
                                                  partyShare seats) ∧
                                      GGRS26CombattingGerrymanderingRCV.paper_seat_share_rounded
                                        (GGRS26CombattingGerrymanderingRCV.paper_pav_selected_seat_count partyShare
                                          seats)
                                        partyShare seats
\end{ReviewVerbatim}

\paragraph{Retained governing declarations}
\begin{itemize}\raggedright
\item \hyperlink{governing-declaration-4c0320b08e4fb484}{Applied\allowbreak{}Modeling\allowbreak{}Lib.\allowbreak{}Social\allowbreak{}Choice.\allowbreak{}Voting.\allowbreak{}Ballot}
\item \hyperlink{governing-declaration-d335879a7207e841}{Applied\allowbreak{}Modeling\allowbreak{}Lib.\allowbreak{}Social\allowbreak{}Choice.\allowbreak{}Voting.\allowbreak{}STV\allowbreak{}Quota}
\item \hyperlink{governing-declaration-3bfffc20e8b6444f}{GGRS26\allowbreak{}Combatting\allowbreak{}Gerrymandering\allowbreak{}RCV.\allowbreak{}Ballot\allowbreak{}Routed\allowbreak{}D\allowbreak{}Favoring\allowbreak{}STV\allowbreak{}Run}
\item \hyperlink{governing-declaration-8b872df78d871e1b}{GGRS26\allowbreak{}Combatting\allowbreak{}Gerrymandering\allowbreak{}RCV.\allowbreak{}Ballot\allowbreak{}Routed\allowbreak{}STV\allowbreak{}State}
\item \hyperlink{governing-declaration-27b5b75da6f75108}{GGRS26\allowbreak{}Combatting\allowbreak{}Gerrymandering\allowbreak{}RCV.\allowbreak{}Ballot\allowbreak{}Routed\allowbreak{}STV\allowbreak{}Terminal}
\item \hyperlink{governing-declaration-9554b7918969e197}{GGRS26\allowbreak{}Combatting\allowbreak{}Gerrymandering\allowbreak{}RCV.\allowbreak{}Ballot\allowbreak{}Routed\allowbreak{}STV\allowbreak{}Transfer\allowbreak{}Policy}
\item \hyperlink{governing-declaration-21463db3a6da8b5e}{GGRS26\allowbreak{}Combatting\allowbreak{}Gerrymandering\allowbreak{}RCV.\allowbreak{}ballot\allowbreak{}Routed\allowbreak{}Party\allowbreak{}Final\allowbreak{}Seats}
\item \hyperlink{governing-declaration-ab89c07c272f0ccf}{GGRS26\allowbreak{}Combatting\allowbreak{}Gerrymandering\allowbreak{}RCV.\allowbreak{}paper\_\allowbreak{}complete\_\allowbreak{}party\_\allowbreak{}ranking\_\allowbreak{}ballots}
\item \hyperlink{paper-prerequisite-c0f468d3ce4de7e0}{GGRS26\allowbreak{}Combatting\allowbreak{}Gerrymandering\allowbreak{}RCV.\allowbreak{}paper\_\allowbreak{}pav\_\allowbreak{}selected\_\allowbreak{}seat\_\allowbreak{}count}
\item \hyperlink{governing-declaration-a26378e50367c86c}{GGRS26\allowbreak{}Combatting\allowbreak{}Gerrymandering\allowbreak{}RCV.\allowbreak{}paper\_\allowbreak{}proposition1\_\allowbreak{}ballot\_\allowbreak{}routed\_\allowbreak{}stv\_\allowbreak{}initial\_\allowbreak{}state}
\item \hyperlink{governing-declaration-64adc81e0c40ba7a}{GGRS26\allowbreak{}Combatting\allowbreak{}Gerrymandering\allowbreak{}RCV.\allowbreak{}paper\_\allowbreak{}seat\_\allowbreak{}share\_\allowbreak{}rounded}
\end{itemize}
\paragraph{Lean-elaborated claim roles}\begin{ReviewVerbatim}
1. parameter: Type u_1
2. parameter: Type u_2
3. parameter: Finset Voter
4. parameter: Finset Voter
5. parameter: Finset Voter
6. parameter: Voter → AppliedModelingLib.SocialChoice.Voting.Ballot Candidate
7. parameter: Finset Candidate
8. parameter: Finset Candidate
9. parameter: Finset Candidate
10. parameter: ℕ
11. parameter: ℕ
12. parameter: ℝ
13. assumption: 1 ≤ seats
14. assumption: 0 ≤ partyShare
15. assumption: partyShare ≤ 1
16. assumption: seats * (seats + 1) ≤ voters
17. assumption: seats ≤ partyCandidates.card
18. assumption: seats ≤ otherPartyCandidates.card
19. assumption: GGRS26CombattingGerrymanderingRCV.paper_complete_party_ranking_ballots partyVoters ballots partyCandidates initialActive
20. assumption: GGRS26CombattingGerrymanderingRCV.paper_complete_party_ranking_ballots otherPartyVoters ballots otherPartyCandidates
  initialActive
21. assumption: voters = allVoters.card
22. assumption: allVoters = partyVoters ∪ otherPartyVoters
23. assumption: Disjoint partyVoters otherPartyVoters
24. assumption: Disjoint partyCandidates otherPartyCandidates
25. assumption: partyCandidates ⊆ initialActive
26. assumption: otherPartyCandidates ⊆ initialActive
27. assumption: initialActive ⊆ partyCandidates ∪ otherPartyCandidates
28. assumption: partyShare * ↑voters = ↑partyVoters.card
29. assumption: (1 - partyShare) * ↑voters = ↑otherPartyVoters.card
30. conclusion: (∀
    (policy :
      GGRS26CombattingGerrymanderingRCV.BallotRoutedSTVTransferPolicy allVoters ballots
        ↑(AppliedModelingLib.SocialChoice.Voting.STVQuota seats voters)),
    (∃ (terminal : GGRS26CombattingGerrymanderingRCV.BallotRoutedSTVState allVoters initialActive),
        GGRS26CombattingGerrymanderingRCV.BallotRoutedDFavoringSTVRun ballots
            (↑(AppliedModelingLib.SocialChoice.Voting.STVQuota seats voters)) policy seats
            (GGRS26CombattingGerrymanderingRCV.paper_proposition1_ballot_routed_stv_initial_state allVoters
              initialActive)
            terminal ∧
          GGRS26CombattingGerrymanderingRCV.BallotRoutedSTVTerminal seats terminal) ∧
      ∀ (terminal : GGRS26CombattingGerrymanderingRCV.BallotRoutedSTVState allVoters initialActive),
        GGRS26CombattingGerrymanderingRCV.BallotRoutedDFavoringSTVRun ballots
            (↑(AppliedModelingLib.SocialChoice.Voting.STVQuota seats voters)) policy seats
            (GGRS26CombattingGerrymanderingRCV.paper_proposition1_ballot_routed_stv_initial_state allVoters
              initialActive)
            terminal →
          GGRS26CombattingGerrymanderingRCV.BallotRoutedSTVTerminal seats terminal →
            GGRS26CombattingGerrymanderingRCV.paper_seat_share_rounded
              (GGRS26CombattingGerrymanderingRCV.ballotRoutedPartyFinalSeats partyCandidates seats terminal) partyShare
              seats) ∧
  GGRS26CombattingGerrymanderingRCV.paper_seat_share_rounded
    (GGRS26CombattingGerrymanderingRCV.paper_pav_selected_seat_count partyShare seats) partyShare seats
\end{ReviewVerbatim}

\paragraph{Lean proof endpoint (built separately)}\begin{ReviewVerbatim}
GGRS26CombattingGerrymanderingRCV.paper_proposition1_source_selected_stv_and_pav
\end{ReviewVerbatim}

\paragraph{Recorded source-to-Spec screening}
\textbf{Verdict:} \texttt{matches}
\\
{\raggedright
\textbf{Reason:} The target retains V >= M(M+1), the two-\allowbreak{}party solid-\allowbreak{}coalition model, at least M candidates per party, exact Droop quota, D-\allowbreak{}favoring STV and the floor-\allowbreak{}or-\allowbreak{}ceiling STV/\allowbreak{}PAV conclusion;\allowbreak{} within-\allowbreak{}party randomness is used only as pathwise outcome support, not as a probability-\allowbreak{}law claim.\allowbreak{}
\par}
\reviewmatch{Matches source input}
\noindent\textbf{Reviewer annotation}\par
\reviewerbox

\end{document}
